\chapter{Unit 4: Quantum Mechanics and Wave-Particle Duality}

\section{Section 1: 100 Chapter-Wise Multiple-Choice Questions (MCQs)}

\subsection*{Part I: Foundational Level (Questions 1 to 50)}

\mcqitem{1}{According to the de Broglie hypothesis, the matter wavelength $\lambda$ of a particle with linear momentum $p$ is:}{$\lambda = \frac{h}{p}$}{$\lambda = \frac{p}{h}$}{$\lambda = h p$}{$\lambda = \frac{\hbar}{p^2}$}
\mcqitem{2}{The de Broglie wavelength of an electron accelerated from rest through an electrostatic potential difference $V$ is given by:}{$\lambda = \frac{1.227}{\sqrt{V}}\text{ nm}$}{$\lambda = \frac{12.27}{\sqrt{V}}\text{ nm}$}{$\lambda = 1.227\sqrt{V}\text{ nm}$}{$\lambda = \frac{0.1227}{V}\text{ nm}$}
\mcqitem{3}{The wave nature of electrons was first directly validated experimentally by:}{The Michelson-Morley experiment}{The Davisson-Germer electron diffraction experiment}{The Stern-Gerlach experiment}{The Millikan oil drop experiment}
\mcqitem{4}{In the Davisson-Germer experiment, the primary diffraction peak was observed at an accelerating potential and scattering angle of:}{$54\text{ V}$ and $\phi = 50^\circ$}{$100\text{ V}$ and $\phi = 90^\circ$}{$25\text{ V}$ and $\phi = 30^\circ$}{$500\text{ V}$ and $\phi = 45^\circ$}
\mcqitem{5}{The phase velocity $v_p$ of a harmonic wave is defined as:}{$v_p = \frac{\omega}{k}$}{$v_p = \frac{d\omega}{dk}$}{$v_p = \omega k$}{$v_p = \frac{k}{\omega}$}
\mcqitem{6}{The group velocity $v_g$ of a wave packet is defined as:}{$v_g = \frac{d\omega}{dk}$}{$v_g = \frac{\omega}{k}$}{$v_g = \frac{k}{\omega}$}{$v_g = \frac{d^2\omega}{dk^2}$}
\mcqitem{7}{For a non-relativistic free particle of mass $m$ and velocity $v$, the group velocity of its de Broglie wave packet equals:}{$v_g = v$ (the physical particle velocity)}{$v_g = v/2$}{$v_g = 2v$}{$v_g = c$}
\mcqitem{8}{For relativistic de Broglie matter waves, the product of phase velocity $v_p$ and group velocity $v_g$ satisfies:}{$v_p \cdot v_g = c^2$}{$v_p \cdot v_g = 1$}{$v_p / v_g = c$}{$v_p + v_g = c$}
\mcqitem{9}{Heisenberg's Uncertainty Principle for position and linear momentum states:}{$\Delta x \cdot \Delta p_x \ge \frac{\hbar}{2}$}{$\Delta x \cdot \Delta p_x = 0$}{$\Delta x \cdot \Delta p_x \le \hbar$}{$\frac{\Delta x}{\Delta p_x} \ge \hbar$}
\mcqitem{10}{Heisenberg's Uncertainty Principle for energy and time is expressed as:}{$\Delta E \cdot \Delta t \ge \frac{\hbar}{2}$}{$\Delta E \cdot \Delta t \le \hbar$}{$\Delta E = \Delta t$}{$\Delta E \cdot \Delta t = 0$}
\mcqitem{11}{Why cannot electrons exist inside the atomic nucleus ($R \sim 10^{-14}\text{ m}$)?}{Uncertainty in momentum would require electron kinetic energy $> 20\text{ MeV}$, which far exceeds observed $\beta$-decay energies ($\sim 2\text{--}3\text{ MeV}$)}{Electrons have positive mass}{Nuclear forces repel negative charges}{Electrons travel at speed $2c$}
\mcqitem{12}{Max Born's statistical interpretation of the complex quantum wavefunction $\Psi(\vec{r},t)$ states that:}{$|\Psi(\vec{r},t)|^2 = \Psi^* \Psi$ represents the probability density of finding the particle at position $\vec{r}$ at time $t$}{$\Psi$ represents the physical mass density of the electron}{$\Psi$ is the classical electric potential}{$|\Psi|^2$ is the total mechanical energy}
\mcqitem{13}{The global normalization condition for a physically acceptable 1D wave function $\psi(x)$ is:}{$\int_{-\infty}^{+\infty} |\psi(x)|^2 dx = 1$}{$\int_{-\infty}^{+\infty} \psi(x) dx = 0$}{$\int_{-\infty}^{+\infty} |\psi(x)|^2 dx = \infty$}{$\psi(0) = 1$}
\mcqitem{14}{Which of the following is NOT a required mathematical condition for a physically well-behaved wave function $\psi(x)$?}{$\psi(x)$ must be single-valued everywhere}{$\psi(x)$ must be continuous everywhere}{$\frac{d\psi}{dx}$ must be continuous everywhere (for finite potentials)}{$\psi(x)$ must become infinite at the boundaries}
\mcqitem{15}{The quantum mechanical operator for linear momentum $\hat{p}_x$ in coordinate space is:}{$\hat{p}_x = -i\hbar \frac{\partial}{\partial x}$}{$\hat{p}_x = i\hbar \frac{\partial}{\partial x}$}{$\hat{p}_x = -\frac{\hbar^2}{2m}\frac{\partial^2}{\partial x^2}$}{$\hat{p}_x = \hbar \frac{\partial}{\partial t}$}
\mcqitem{16}{The quantum mechanical total energy operator $\hat{E}$ is given by:}{$\hat{E} = i\hbar \frac{\partial}{\partial t}$}{$\hat{E} = -i\hbar \frac{\partial}{\partial t}$}{$\hat{E} = \hbar \frac{\partial}{\partial x}$}{$\hat{E} = -\frac{\hbar^2}{2m}$}
\mcqitem{17}{The Hamiltonian operator $\hat{H}$ for a particle of mass $m$ in a potential $V(x)$ is:}{$\hat{H} = -\frac{\hbar^2}{2m}\frac{d^2}{dx^2} + V(x)$}{$\hat{H} = \frac{\hbar^2}{2m}\frac{d^2}{dx^2} + V(x)$}{$\hat{H} = -i\hbar \frac{d}{dx} + V(x)$}{$\hat{H} = -\frac{\hbar^2}{2m}\nabla^2$}
\mcqitem{18}{The 1D Time-Dependent Schrödinger Equation (TDSE) is written as:}{$i\hbar \frac{\partial \Psi}{\partial t} = -\frac{\hbar^2}{2m}\frac{\partial^2 \Psi}{\partial x^2} + V(x)\Psi$}{$i\hbar \frac{\partial \Psi}{\partial x} = -\frac{\hbar^2}{2m}\frac{\partial^2 \Psi}{\partial t^2} + V\Psi$}{$\frac{\partial^2 \Psi}{\partial t^2} = c^2 \frac{\partial^2 \Psi}{\partial x^2}$}{$\hat{H}\Psi = 0$}
\mcqitem{19}{The 1D Time-Independent Schrödinger Equation (TISE) for stationary states of energy $E$ is:}{$-\frac{\hbar^2}{2m}\frac{d^2 \psi}{dx^2} + V(x)\psi = E\psi$}{$\frac{d^2\psi}{dx^2} + \frac{2m}{\hbar^2}(E - V)\psi = 0$}{$\hat{H}\psi = E\psi$}{All of the above are equivalent forms}
\mcqitem{20}{For a particle of mass $m$ trapped in a 1D infinite square well of width $L$ ($0 \le x \le L$), the quantized energy eigenvalues $E_n$ are:}{$E_n = \frac{n^2 h^2}{8mL^2} = \frac{n^2 \pi^2 \hbar^2}{2mL^2} \quad (n = 1, 2, 3,\dots)$}{$E_n = \frac{n h}{8mL}$}{$E_n = \frac{n^2 \hbar^2}{2mL}$}{$E_n = \frac{(n+1/2)\hbar^2}{2mL^2}$}
\mcqitem{21}{Why is the quantum number $n=0$ physically impossible for a particle in an infinite potential well?}{Because $n=0$ gives $\psi(x) = \sqrt{2/L}\sin(0) = 0$ everywhere, which means no particle exists (violates normalization)}{Because energy would be infinite}{Because momentum would be negative}{Because the mass becomes zero}
\mcqitem{22}{The zero-point energy (ground-state energy $E_1$) of a particle in an infinite square well is:}{$E_1 = 0$}{$E_1 = \frac{h^2}{8mL^2} > 0$}{$E_1 = \infty$}{$E_1 = -\frac{h^2}{8mL^2}$}
\mcqitem{23}{The normalized spatial wave functions $\psi_n(x)$ for a particle in a 1D infinite well of width $L$ ($0 \le x \le L$) are:}{$\psi_n(x) = \sqrt{\frac{2}{L}}\sin\left(\frac{n\pi x}{L}\right)$}{$\psi_n(x) = \sqrt{\frac{1}{L}}\cos\left(\frac{n\pi x}{L}\right)$}{$\psi_n(x) = \frac{2}{L}\sin\left(\frac{n\pi x}{L}\right)$}{$\psi_n(x) = \sqrt{\frac{2}{L}}e^{i n\pi x/L}$}
\mcqitem{24}{The number of interior spatial nodes (points where $\psi(x) = 0$ inside the well) for the $n$-th energy state is:}{$n$}{$n - 1$}{$n + 1$}{$2n$}
\mcqitem{25}{The expectation value $\langle x \rangle$ of position for any stationary state $\psi_n(x)$ in a 1D box of width $L$ ($0 \le x \le L$) is:}{$\langle x \rangle = 0$}{$\langle x \rangle = \frac{L}{2}$}{$\langle x \rangle = \frac{L}{4}$}{$\langle x \rangle = L$}
\mcqitem{26}{The expectation value of linear momentum $\langle p_x \rangle$ for any stationary state in a 1D infinite well is:}{$\langle p_x \rangle = 0$}{$\langle p_x \rangle = \frac{n\pi\hbar}{L}$}{$\langle p_x \rangle = \frac{n h}{2L}$}{$\langle p_x \rangle = \infty$}
\mcqitem{27}{The ratio of the first four energy levels ($E_1 : E_2 : E_3 : E_4$) for a particle in a 1D infinite potential well is:}{$1 : 2 : 3 : 4$}{$1 : 4 : 9 : 16$}{$1 : 3 : 5 : 7$}{$1 : 8 : 27 : 64$}
\mcqitem{28}{For a particle in a 3D cubical box of side $L$, the energy eigenvalues are given by:}{$E_{n_x,n_y,n_z} = \frac{h^2}{8mL^2}(n_x^2 + n_y^2 + n_z^2)$}{$E = \frac{h^2}{8mL^2}(n_x + n_y + n_z)^2$}{$E = \frac{3 h^2}{8mL^2} n$}{$E = \frac{h^2}{8mL^2}(n_x n_y n_z)$}
\mcqitem{29}{In a 3D cubical box, the degree of degeneracy of the first excited energy level ($n_x^2 + n_y^2 + n_z^2 = 1^2 + 1^2 + 2^2 = 6$) is:}{Non-degenerate ($1$)}{$2$-fold degenerate}{$3$-fold degenerate ($(2,1,1), (1,2,1), (1,1,2)$)}{$6$-fold degenerate}
\mcqitem{30}{Quantum mechanical tunneling is the physical phenomenon where a particle:}{Penetrates through a potential energy barrier whose height $V_0$ is greater than the particle's total energy $E$}{Gains infinite velocity}{Destroys the potential barrier mechanically}{Absorbs an infinite number of photons}
\mcqitem{31}{Which prominent scientific instrument operates directly on the principle of quantum mechanical electron tunneling?}{Scanning Tunneling Microscope (STM)}{Optical telescope}{Cathode ray oscilloscope}{Geiger-Muller counter}
\mcqitem{32}{The transmission probability $T$ for a particle of energy $E$ tunneling through a rectangular barrier of height $V_0 > E$ and width $a$ scales as:}{$T \propto e^{-2\kappa a}$ where $\kappa = \sqrt{\frac{2m(V_0 - E)}{\hbar^2}}$}{$T \propto e^{+\kappa a}$}{$T \propto 1/a$}{$T \propto (V_0 - E)^2$}
\mcqitem{33}{The commutator bracket $[\hat{x}, \hat{p}_x] = \hat{x}\hat{p}_x - \hat{p}_x\hat{x}$ evaluates to:}{$0$}{$i\hbar$}{$-i\hbar$}{$\hbar^2$}
\mcqitem{34}{The non-commutativity of operators $[\hat{A}, \hat{B}] \ne 0$ fundamentally implies that:}{The physical observables $A$ and $B$ cannot be simultaneously measured with arbitrary precision}{Both observables must equal zero}{The system has zero energy}{The wave function is unnormalizable}
\mcqitem{35}{Hermitian operators in quantum mechanics are mathematically essential because:}{Their eigenvalues are guaranteed to be strictly real numbers, matching physical experimental measurements}{They eliminate imaginary numbers from equations}{They are always diagonal matrices}{They commute with all other operators}
\mcqitem{36}{Two wavefunctions $\psi_m(x)$ and $\psi_n(x)$ corresponding to distinct energy eigenvalues of a Hermitian Hamiltonian are:}{Mutually orthogonal: $\int \psi_m^*(x)\psi_n(x)dx = 0 \quad (m \ne n)$}{Identical}{Linearly dependent}{Infinite at the origin}
\mcqitem{37}{The expectation value $\langle A \rangle$ of a physical observable associated with operator $\hat{A}$ for state $\Psi$ is calculated as:}{$\langle A \rangle = \int_{-\infty}^{+\infty} \Psi^* \hat{A} \Psi\,dx$}{$\langle A \rangle = \int_{-\infty}^{+\infty} \hat{A} |\Psi|^2\,dx$}{$\langle A \rangle = \frac{\hat{A}}{\int |\Psi|^2 dx}$}{$\langle A \rangle = \Psi^* \hat{A}$}
\mcqitem{38}{In Compton scattering, the shift in photon wavelength $\Delta \lambda = \lambda' - \lambda$ after scattering by angle $\theta$ is:}{$\Delta \lambda = \frac{h}{m_e c}(1 - \cos\theta)$}{$\Delta \lambda = \frac{h}{m_e c}(1 + \cos\theta)$}{$\Delta \lambda = \frac{m_e c}{h}(1 - \sin\theta)$}{$\Delta \lambda = \frac{h}{2m_e c}\sin\theta$}
\mcqitem{39}{The Compton wavelength of an electron $\lambda_C = \frac{h}{m_e c}$ has the numerical value:}{$0.0243\text{ \AA} = 2.43\text{ pm}$}{$0.243\text{ nm}$}{$2.43\text{ nm}$}{$24.3\text{ nm}$}
\mcqitem{40}{The maximum possible Compton wavelength shift $\Delta \lambda_{\text{max}}$ occurs at a scattering angle of:}{$\theta = 180^\circ$ (backscattering), where $\Delta\lambda = 2\lambda_C \approx 0.0485\text{ \AA}$}{$\theta = 90^\circ$}{$\theta = 0^\circ$}{$\theta = 45^\circ$}
\mcqitem{41}{The photoelectric work function $\Phi$ of a metal represents:}{The minimum photon energy required to liberate an electron from the metal surface at $T=0\text{ K}$}{The maximum kinetic energy of photoelectrons}{The energy needed to melt the metal}{The total electrostatic potential of the lattice}
\mcqitem{42}{Einstein's photoelectric equation is given by:}{$K_{\text{max}} = h\nu - \Phi = e V_0$}{$K_{\text{max}} = h\nu + \Phi$}{$K_{\text{max}} = \frac{h\nu}{\Phi}$}{$K_{\text{max}} = \Phi - h\nu$}
\mcqitem{43}{The slope of the graph of stopping potential $V_0$ versus frequency $\nu$ in a photoelectric experiment is equal to:}{$h/e$}{$e/h$}{$h$}{$e$}
\mcqitem{44}{An electron is accelerated through a potential difference $V = 100\text{ V}$. Its de Broglie wavelength is:}{$0.123\text{ nm} = 1.23\text{ \AA}$}{$1.23\text{ nm}$}{$0.0123\text{ nm}$}{$12.3\text{ nm}$}
\mcqitem{45}{A thermal neutron at temperature $T = 300\text{ K}$ has average kinetic energy $E_k = \frac{3}{2}k_B T$. Its de Broglie wavelength is approximately:}{$0.145\text{ nm} = 1.45\text{ \AA}$}{$1.45\text{ nm}$}{$0.0145\text{ nm}$}{$14.5\text{ nm}$}
\mcqitem{46}{An electron is confined in a 1D box of width $L = 0.1\text{ nm}$ ($1\text{ \AA}$). Its ground-state energy $E_1$ is approximately:}{$37.6\text{ eV}$}{$3.76\text{ eV}$}{$376\text{ eV}$}{$0.376\text{ eV}$}
\mcqitem{47}{For the electron in Question 46, the first excited state energy $E_2$ is:}{$150.4\text{ eV}$}{$75.2\text{ eV}$}{$112.8\text{ eV}$}{$338.4\text{ eV}$}
\mcqitem{48}{The probability of finding a particle in the ground state of an infinite well ($0 \le x \le L$) in the middle third of the box $[L/3, 2L/3]$ is:}{$\frac{1}{3} + \frac{\sqrt{3}}{2\pi} \approx 60.9\%$}{$33.3\%$}{$50.0\%$}{$25.0\%$}
\mcqitem{49}{If the position of an electron is measured with an uncertainty $\Delta x = 1.0\text{ \AA} = 10^{-10}\text{ m}$, the minimum uncertainty in its velocity $\Delta v_x$ is approximately:}{$5.79 \times 10^5\text{ m/s}$}{$5.79 \times 10^3\text{ m/s}$}{$1.16 \times 10^6\text{ m/s}$}{$3.00 \times 10^8\text{ m/s}$}
\mcqitem{50}{Ehrenfest's Theorem in quantum mechanics proves that:}{Quantum mechanical expectation values follow classical Newtonian laws of motion in the correspondence limit ($\frac{d\langle x \rangle}{dt} = \frac{\langle p \rangle}{m}$ and $\frac{d\langle p \rangle}{dt} = -\langle \nabla V \rangle$)}{Quantum mechanics replaces all conservation laws}{Energy is never conserved in quantum systems}{Electrons travel on fixed Keplerian orbits}

\subsection*{Part II: Intermediate Analytical Level (Questions 51 to 80)}

\mcqitem{51}{A macroscopic marble of mass $m = 10\text{ g}$ moves at $v = 10\text{ m/s}$. Its de Broglie wavelength is:}{$6.63 \times 10^{-33}\text{ m}$}{$6.63 \times 10^{-30}\text{ m}$}{$6.63 \times 10^{-27}\text{ m}$}{$1.00 \times 10^{-15}\text{ m}$}
\mcqitem{52}{The kinetic energy of an electron whose de Broglie wavelength equals its Compton wavelength ($\lambda = \lambda_C = h/(m_e c)$) is (relativistic calculation):}{$(\sqrt{2} - 1)m_e c^2 \approx 0.414 m_e c^2 \approx 212\text{ keV}$}{$0.5 m_e c^2$}{$1.0 m_e c^2$}{$2.0 m_e c^2$}
\mcqitem{53}{An excited atomic energy state has an average radiative lifetime $\tau = 1.0\text{ ns} = 10^{-9}\text{ s}$. The minimum natural energy linewidth $\Delta E$ of the emitted line is:}{$3.29 \times 10^{-7}\text{ eV}$}{$3.29 \times 10^{-5}\text{ eV}$}{$6.58 \times 10^{-16}\text{ eV}$}{$1.00\text{ eV}$}
\mcqitem{54}{For the atomic state in Question 53 emitting at $\lambda_0 = 500\text{ nm}$ ($\nu_0 = 6 \times 10^{14}\text{ Hz}$), the fractional natural line broadening $\Delta \nu / \nu_0$ is:}{$1.33 \times 10^{-7}$}{$1.33 \times 10^{-5}$}{$1.33 \times 10^{-9}$}{$5.00 \times 10^{-3}$}
\mcqitem{55}{A proton and an electron have the same kinetic energy. The ratio of their de Broglie wavelengths $\lambda_p / \lambda_e$ is:}{$\sqrt{\frac{m_e}{m_p}} \approx \frac{1}{\sqrt{1836}} \approx \frac{1}{42.85}$}{$\sqrt{\frac{m_p}{m_e}} \approx 42.85$}{$\frac{m_e}{m_p} \approx \frac{1}{1836}$}{$1.0$}
\mcqitem{56}{A proton and an electron have the exact same momentum $p$. The ratio of their de Broglie wavelengths $\lambda_p / \lambda_e$ is:}{$1.0$}{$1836$}{$1/1836$}{$42.85$}
\mcqitem{57}{For an infinite square well of width $L$, the wavelength of a photon emitted during the transition from $n=3$ to $n=1$ is:}{$\lambda = \frac{hc}{8 E_1}$}{$\lambda = \frac{hc}{4 E_1}$}{$\lambda = \frac{hc}{9 E_1}$}{$\lambda = \frac{8 E_1}{hc}$}
\mcqitem{58}{The variance in position $\sigma_x^2 = \langle x^2 \rangle - \langle x \rangle^2$ for the ground state $n=1$ in a 1D box of width $L$ is:}{$L^2\left(\frac{1}{12} - \frac{1}{2\pi^2}\right) \approx 0.0326 L^2$}{$\frac{L^2}{12}$}{$\frac{L^2}{4}$}{$0$}
\mcqitem{59}{The variance in momentum $\sigma_p^2 = \langle p^2 \rangle - \langle p \rangle^2$ for the $n=1$ ground state in the box is:}{$\frac{\pi^2 \hbar^2}{L^2}$}{$\frac{\pi \hbar}{L}$}{$0$}{$\frac{h^2}{4L^2}$}
\mcqitem{60}{The uncertainty product $\Delta x \cdot \Delta p$ for the ground state of a particle in an infinite well evaluates to:}{$\hbar \sqrt{\frac{\pi^2}{12} - \frac{1}{2}} \approx 0.568 \hbar > \frac{\hbar}{2}$}{$\frac{\hbar}{2}$}{$0$}{$\pi \hbar$}
\mcqitem{61}{The quantum mechanical probability current density $\vec{J}(\vec{r},t)$ is defined as:}{$\vec{J} = \frac{\hbar}{2mi}\left(\Psi^* \nabla \Psi - \Psi \nabla \Psi^*\right)$}{$\vec{J} = |\Psi|^2 \vec{v}$}{$\vec{J} = \Psi^* \Psi$}{$\vec{J} = -\frac{\hbar^2}{2m}\nabla \Psi$}
\mcqitem{62}{For a plane wave state $\Psi(x,t) = A e^{i(kx - \omega t)}$, the probability current density $J_x$ is:}{$J_x = |A|^2 \frac{\hbar k}{m} = |A|^2 v$}{$J_x = 0$}{$J_x = |A|^2$}{$J_x = \hbar k$}
\mcqitem{63}{For any real bound-state spatial eigenfunction $\psi(x) = \psi^*(x)$ (standing wave), the net probability current density $J_x$ is:}{Strictly zero ($J_x = 0$ everywhere)}{$1.0$}{$\infty$}{Proportional to $x$}
\mcqitem{64}{An electron has total energy $E = 2.0\text{ eV}$ incident on a finite potential step of height $V_0 = 5.0\text{ eV}$. In the barrier region ($x > 0$), the wave function behaves as:}{A pure decaying exponential $\psi(x) = C e^{-\kappa x}$ (evanescent wave)}{A purely propagating sine wave}{Zero identically everywhere}{An exponentially exploding wave $e^{+\kappa x}$}
\mcqitem{65}{For the electron in Question 64, the penetration decay constant $\kappa$ is approximately:}{$8.87 \times 10^9\text{ m}^{-1} = 0.887\text{ \AA}^{-1}$}{$1.00 \times 10^7\text{ m}^{-1}$}{$4.50 \times 10^5\text{ m}^{-1}$}{$2.20 \times 10^9\text{ m}^{-1}$}
\mcqitem{66}{The characteristic penetration depth $d = 1/\kappa$ over which the electron wave amplitude drops to $1/e \approx 36.8\%$ is:}{$1.13\text{ \AA} = 0.113\text{ nm}$}{$1.13\text{ nm}$}{$11.3\text{ nm}$}{$0.011\text{ nm}$}
\mcqitem{67}{The reflection coefficient $R$ for the potential step in Question 64 ($E < V_0$) is:}{$R = 1.0$ ($100\%$ total quantum reflection)}{$R = 0.5$}{$R = 0$}{$R = 0.25$}
\mcqitem{68}{A particle in a 1D harmonic oscillator potential $V(x) = \frac{1}{2}m\omega^2 x^2$ has quantized energy levels:}{$E_n = (n + 1/2)\hbar\omega \quad (n = 0, 1, 2,\dots)$}{$E_n = n \hbar\omega$}{$E_n = n^2 \hbar\omega$}{$E_n = (n + 1)\hbar\omega$}
\mcqitem{69}{The zero-point energy of a 1D quantum harmonic oscillator of frequency $\omega$ is:}{$E_0 = \frac{1}{2}\hbar\omega$}{$E_0 = 0$}{$E_0 = \hbar\omega$}{$E_0 = \frac{3}{2}\hbar\omega$}
\mcqitem{70}{In the Bohr model of the hydrogen atom, the orbital angular momentum $L$ of the electron is quantized according to:}{$L = n \hbar = \frac{n h}{2\pi} \quad (n = 1, 2, 3,\dots)$}{$L = n h$}{$L = \frac{\hbar}{n}$}{$L = n^2 \hbar$}
\mcqitem{71}{Bohr's angular momentum quantization $m v r = n\hbar$ is physically explained by de Broglie matter waves as:}{The requirement that the circumference of the electron orbit must fit an exact integer number of de Broglie wavelengths ($2\pi r = n\lambda$)}{Total internal reflection inside the proton}{Relativistic mass increase}{Heisenberg momentum cutoff}
\mcqitem{72}{The radius of the first Bohr orbit $a_0$ in hydrogen is given by:}{$a_0 = \frac{4\pi \epsilon_0 \hbar^2}{m_e e^2} \approx 0.529\text{ \AA} = 0.0529\text{ nm}$}{$a_0 = 0.0529\text{ \AA}$}{$a_0 = 5.29\text{ nm}$}{$a_0 = 1.00\text{ \AA}$}
\mcqitem{73}{The ionization energy of hydrogen in its ground state ($n=1$) is:}{$13.6\text{ eV}$}{$3.4\text{ eV}$}{$1.51\text{ eV}$}{$54.4\text{ eV}$}
\mcqitem{74}{The Rydberg constant $R_\infty = \frac{m_e e^4}{8 \epsilon_0^2 h^3 c}$ has the approximate value:}{$1.097 \times 10^7\text{ m}^{-1}$}{$1.097 \times 10^5\text{ m}^{-1}$}{$3.29 \times 10^{15}\text{ Hz}$}{$6.626 \times 10^{-34}\text{ J}\cdot\text{s}$}
\mcqitem{75}{The Lyman spectral series of hydrogen occurs in which region of the electromagnetic spectrum?}{Ultraviolet (UV)}{Visible}{Infrared (IR)}{Far-Infrared}
\mcqitem{76}{The Balmer spectral series of hydrogen ($n \to 2$) is unique because:}{Its principal emission lines ($H_\alpha, H_\beta, H_\gamma, H_\delta$) lie directly in the visible light spectrum ($400\text{--}700\text{ nm}$)}{It emits X-rays only}{It has zero transition probability}{It requires nuclear fission}
\mcqitem{77}{If a quantum state is a linear superposition $\Psi(x) = c_1 \psi_1(x) + c_2 \psi_2(x)$, the probability of measuring energy $E_1$ is:}{$|c_1|^2$}{$|c_1|$}{$c_1 + c_2$}{$|c_1|^2 / |c_2|^2$}
\mcqitem{78}{The parity of the wave function $\psi_n(x)$ in a symmetric potential $V(-x) = V(x)$ is:}{Even for $n=1, 3, 5\dots$ (odd $n$) and odd for $n=2, 4, 6\dots$ (even $n$)}{Always positive}{Always negative}{Undefined}
\mcqitem{79}{The degeneracy of the energy level $E = 14 E_1$ for a particle in a 3D cubical box ($n_x^2 + n_y^2 + n_z^2 = 14$) is:}{$6$-fold degenerate}{$3$-fold degenerate}{$1$-fold (non-degenerate)}{$12$-fold degenerate}
\mcqitem{80}{The expectation value $\langle p_x^2 \rangle$ for state $\psi_n(x)$ in a 1D box of width $L$ is:}{$\frac{n^2 \pi^2 \hbar^2}{L^2}$}{$0$}{$\frac{n \pi \hbar}{L}$}{$\frac{n^2 h^2}{4L^2}$}

\subsection*{Part III: Advanced \& Numerical Problems (Questions 81 to 100)}

\mcqitem{81}{The time evolution of a non-stationary state $\Psi(x,t) = \frac{1}{\sqrt{2}}\psi_1(x)e^{-iE_1 t/\hbar} + \frac{1}{\sqrt{2}}\psi_2(x)e^{-iE_2 t/\hbar}$ produces a probability density $|\Psi(x,t)|^2$ that oscillates in time with angular frequency:}{$\omega_{21} = \frac{E_2 - E_1}{\hbar}$ (Bohr transition frequency)}{$\omega = \frac{E_1 + E_2}{\hbar}$}{$\omega = 0$ (static)}{$\omega = \frac{E_2}{\hbar}$}
\mcqitem{82}{In the WKB (Wentzel-Kramers-Brillouin) semiclassical approximation, the transmission probability $T$ through an arbitrary potential barrier $V(x) > E$ between turning points $x_1$ and $x_2$ is:}{$T \approx \exp\left( -2\int_{x_1}^{x_2} \sqrt{\frac{2m(V(x) - E)}{\hbar^2}} dx \right)$}{$T \approx \exp\left( +\int_{x_1}^{x_2} V(x) dx \right)$}{$T = 0$}{$T = 1 - \frac{V_0}{E}$}
\mcqitem{83}{George Gamow applied the quantum tunneling barrier penetration formula to explain:}{Alpha ($\alpha$) decay lifetimes of radioactive heavy nuclei (Geiger-Nuttall law)}{Beta decay spectrum}{Gamma ray emission}{Cosmic ray creation}
\mcqitem{84}{The Dirac delta function potential $V(x) = -\alpha \delta(x)$ with $\alpha > 0$ supports:}{Exactly one single bound state with energy $E = -\frac{m\alpha^2}{2\hbar^2}$}{Infinite bound states}{Zero bound states}{Two degenerate bound states}
\mcqitem{85}{For a particle in a finite square well of depth $V_0$ and width $2a$, bound states require:}{$\psi(x)$ to oscillate sinusoidally inside the well ($|x| < a$) and decay exponentially outside ($|x| > a$)}{$\psi(x) = 0$ outside the well}{$\psi(x)$ to be infinite at boundaries}{Energy $E > V_0$}
\mcqitem{86}{As the potential barrier height $V_0 \to \infty$, the energy eigenvalues of a finite square well:}{Converge continuously to the eigenvalues of the infinite potential well $E_n = \frac{n^2 h^2}{8mL^2}$}{Drop to zero}{Become negative}{Become continuous}
\mcqitem{87}{The Pauli Exclusion Principle states that:}{Two identical fermions (half-integer spin particles, such as electrons) cannot occupy the exact same quantum state simultaneously}{All bosons must have distinct energies}{Electrons have infinite mass}{Photons cannot share the same mode}
\mcqitem{88}{Particles with integer spin ($s = 0, 1, 2,\dots$) that obey Bose-Einstein statistics and can occupy the same quantum state in unlimited numbers are called:}{Bosons (e.g., photons, gluons, $^4\text{He}$ atoms)}{Fermions}{Quarks}{Leptons}
\mcqitem{89}{The Fermi energy $E_F(0)$ of a 3D degenerate free electron gas at $T=0\text{ K}$ with electron density $n = N/V$ is:}{$E_F = \frac{\hbar^2}{2m}(3\pi^2 n)^{2/3}$}{$E_F = \frac{\hbar^2}{2m}(2\pi^2 n)^{1/2}$}{$E_F = \frac{h^2}{8m} n$}{$E_F = \frac{3}{5} n k_B T$}
\mcqitem{90}{The average energy per electron $\langle E \rangle$ in a 3D free electron gas at $T=0\text{ K}$ in terms of Fermi energy $E_F$ is:}{$\langle E \rangle = \frac{3}{5} E_F$}{$\langle E \rangle = \frac{1}{2} E_F$}{$\langle E \rangle = E_F$}{$\langle E \rangle = \frac{5}{3} E_F$}
\mcqitem{91}{In the Kronig-Penney model of a periodic crystal lattice, energy band gaps appear because:}{Destructive Bragg reflection of electron matter waves occurs at Brillouin zone boundaries ($k = \pm n\pi/a$)}{Electrons collide with nuclei and stop}{Lattice ions absorb all kinetic energy}{The speed of light drops to zero}
\mcqitem{92}{The effective mass $m^*$ of an electron in a crystal energy band $E(k)$ is defined by:}{$m^* = \hbar^2 \left( \frac{d^2 E}{dk^2} \right)^{-1}$}{$m^* = \frac{1}{\hbar}\frac{dE}{dk}$}{$m^* = \hbar^2 \frac{d^2 E}{dk^2}$}{$m^* = \frac{E}{c^2}$}
\mcqitem{93}{Near the top of a valence band where the curvature $\frac{d^2 E}{dk^2} < 0$, the effective mass $m^*$ is:}{Negative ($m^* < 0$), which is physically treated as a positively charged 'hole'}{Infinite}{Zero}{Equal to free electron mass $m_0$}
\mcqitem{94}{The density of states $g(E)$ for a 1D quantum wire scales with energy $E$ as:}{$g_{\text{1D}}(E) \propto E^{-1/2}$}{$g_{\text{2D}}(E) = \text{constant}$}{$g_{\text{3D}}(E) \propto E^{+1/2}$}{All of the above statements are correct}
\mcqitem{95}{In quantum dots (0D nanocrystals), the density of states is represented by:}{A series of discrete delta functions $\sum \delta(E - E_i)$ (atomic-like sharp levels)}{A smooth parabolic curve}{A continuous line}{Zero everywhere}
\mcqitem{96}{A quantum well has width $L = 5.0\text{ nm}$ in GaAs ($m_e^* = 0.067 m_0$). The ground-state subband energy $E_1$ for electrons is approximately:}{$22.5\text{ meV}$}{$2.25\text{ eV}$}{$225\text{ meV}$}{$0.225\text{ meV}$}
\mcqitem{97}{The Klein-Gordon equation is a relativistic wave equation that describes:}{Spin-$0$ relativistic scalar bosons (such as pions and Higgs bosons)}{Spin-$1/2$ electrons}{Photons exclusively}{Classical fluid particles}
\mcqitem{98}{The Dirac equation in relativistic quantum mechanics successfully predicted the existence of:}{Antimatter (the positron $e^+$) and electron intrinsic spin $s = 1/2$}{Dark energy}{Black holes}{Sound waves}
\mcqitem{99}{The Aharonov-Bohm effect proves that in quantum mechanics:}{The electromagnetic vector potential $\vec{A}$ has physical, gauge-invariant reality and alters electron phase even in regions where $\vec{B} = 0$}{Magnetic fields do not exist}{Phase cannot be measured}{Electric charge is not conserved}
\mcqitem{100}{Quantum entanglement describes a multi-particle state where:}{The quantum state of each particle cannot be described independently of the state of the other particles, regardless of spatial separation}{Particles are physically glued together with glue}{Electrons have infinite charge}{Light speeds up to infinity}

\newpage
\section{Section 2: Complete Answer Key with Detailed Rationales}

\subsection*{Master Answer Key (Questions 1 to 25)}

{\small
\noindent\begin{tabularx}{\textwidth}{|c|c|X|}
\hline
\textbf{Q\#} & \textbf{Ans} & \textbf{Technical Rationale \& Calculation} \\
\hline
\textbf{1} & \textbf{(a)} & The de Broglie relation states that matter wavelength is inversely proportional to momentum: $\lambda = h/p$. \\ \hline
\textbf{2} & \textbf{(a)} & $\lambda = \frac{h}{\sqrt{2m_e e V}} = \frac{1.227 \times 10^{-9}}{\sqrt{V}}\text{ m} = \frac{1.227}{\sqrt{V}}\text{ nm}$. \\ \hline
\textbf{3} & \textbf{(b)} & Davisson and Germer observed diffraction peaks of electrons scattered from a nickel crystal target in 1927. \\ \hline
\textbf{4} & \textbf{(a)} & A sharp diffraction peak occurred at $V = 54\text{ V}$ at angle $\phi = 50^\circ$, yielding $\lambda = 0.167\text{ nm}$ matching Bragg's law. \\ \hline
\textbf{5} & \textbf{(a)} & Phase velocity is the speed at which individual monochromatic wave crests travel: $v_p = \omega/k$. \\ \hline
\textbf{6} & \textbf{(a)} & Group velocity is the velocity of the envelope modulating the wave packet: $v_g = d\omega/dk$. \\ \hline
\textbf{7} & \textbf{(a)} & Using $E = \hbar\omega = \frac{\hbar^2 k^2}{2m}$, $v_g = \frac{d\omega}{dk} = \frac{\hbar k}{m} = \frac{p}{m} = v$. \\ \hline
\textbf{8} & \textbf{(a)} & Since $E = \hbar\omega = \gamma m c^2$ and $p = \hbar k = \gamma m v$, $v_p \cdot v_g = (\omega/k)(d\omega/dk) = (c^2/v)(v) = c^2$. \\ \hline
\textbf{9} & \textbf{(a)} & The fundamental quantum uncertainty relation is $\Delta x \cdot \Delta p_x \ge \frac{\hbar}{2}$ where $\hbar = h/(2\pi)$. \\ \hline
\textbf{10} & \textbf{(a)} & The conjugate energy-time uncertainty relation is $\Delta E \cdot \Delta t \ge \frac{\hbar}{2}$. \\ \hline
\textbf{11} & \textbf{(a)} & Confinement to $\Delta x \sim 10^{-14}\text{ m}$ gives $\Delta p \approx 5.3 \times 10^{-21}\text{ kg}\cdot\text{m/s}$, requiring $E_k > 20\text{ MeV}$, proving electrons cannot reside in nuclei. \\ \hline
\textbf{12} & \textbf{(a)} & Born established that $|\Psi(\vec{r},t)|^2 dV$ is the probability of finding the quantum particle in differential volume $dV$. \\ \hline
\textbf{13} & \textbf{(a)} & The total probability of finding the particle somewhere across all space must equal $100\%$ ($1.0$). \\ \hline
\textbf{14} & \textbf{(d)} & A valid wavefunction must remain finite and square-integrable everywhere; it must never blow up to infinity. \\ \hline
\textbf{15} & \textbf{(a)} & In the coordinate representation, linear momentum is represented by the differential operator $\hat{p}_x = -i\hbar \frac{\partial}{\partial x}$. \\ \hline
\textbf{16} & \textbf{(a)} & The quantum mechanical energy operator is $\hat{E} = i\hbar \frac{\partial}{\partial t}$. \\ \hline
\textbf{17} & \textbf{(a)} & The Hamiltonian represents total energy $\hat{H} = \hat{T} + \hat{V} = \frac{\hat{p}^2}{2m} + V(x) = -\frac{\hbar^2}{2m}\frac{d^2}{dx^2} + V(x)$. \\ \hline
\textbf{18} & \textbf{(a)} & The TDSE is $i\hbar \frac{\partial \Psi}{\partial t} = \hat{H}\Psi = -\frac{\hbar^2}{2m}\frac{\partial^2 \Psi}{\partial x^2} + V(x)\Psi$. \\ \hline
\textbf{19} & \textbf{(d)} & All three statements are mathematically equivalent representations of the eigenvalue equation $\hat{H}\psi = E\psi$. \\ \hline
\textbf{20} & \textbf{(a)} & Boundary conditions $\psi(0)=\psi(L)=0$ require $k_n = n\pi/L$, giving $E_n = \frac{\hbar^2 k_n^2}{2m} = \frac{n^2 h^2}{8mL^2}$. \\ \hline
\textbf{21} & \textbf{(a)} & $n=0$ yields a vanishing wavefunction $\psi(x) = 0$, representing the complete absence of a particle. \\ \hline
\textbf{22} & \textbf{(b)} & Zero-point energy is strictly non-zero ($E_1 = \frac{h^2}{8mL^2} > 0$), fully consistent with Heisenberg's uncertainty principle. \\ \hline
\textbf{23} & \textbf{(a)} & Integrating $\int_0^L A^2 \sin^2(n\pi x/L)dx = 1$ yields normalization constant $A = \sqrt{2/L}$. \\ \hline
\textbf{24} & \textbf{(b)} & The $n$-th quantum state has exactly $(n - 1)$ interior nodes (e.g., ground state $n=1$ has $0$ interior nodes). \\ \hline
\textbf{25} & \textbf{(b)} & Because probability density $|\psi_n(x)|^2 = \frac{2}{L}\sin^2(n\pi x/L)$ is symmetric about the center, $\langle x \rangle = L/2$. \\ \hline
\end{tabularx}
}

\newpage

\subsection*{Master Answer Key (Questions 26 to 50)}

{\small
\noindent\begin{tabularx}{\textwidth}{|c|c|X|}
\hline
\textbf{Q\#} & \textbf{Ans} & \textbf{Technical Rationale \& Calculation} \\
\hline
\textbf{26} & \textbf{(a)} & The particle oscillates back and forth with equal probability of $+p$ and $-p$, making the net expectation momentum $\langle p_x \rangle = 0$. \\ \hline
\textbf{27} & \textbf{(b)} & Because $E_n = n^2 E_1$, the energy levels scale as $1^2 : 2^2 : 3^2 : 4^2 = 1 : 4 : 9 : 16$. \\ \hline
\textbf{28} & \textbf{(a)} & Separation of variables in 3D yields $E = \frac{h^2}{8mL^2}(n_x^2 + n_y^2 + n_z^2)$ with quantum numbers $n_x, n_y, n_z \in \{1, 2, 3,\dots\}$. \\ \hline
\textbf{29} & \textbf{(c)} & Three distinct combinations $(2,1,1), (1,2,1),$ and $(1,1,2)$ share the exact same energy $E = 6E_1$. \\ \hline
\textbf{30} & \textbf{(a)} & Because wavefunctions decay exponentially rather than dropping abruptly to zero, there is a finite transmission probability $T > 0$ through barriers $V_0 > E$. \\ \hline
\textbf{31} & \textbf{(a)} & Binnig and Rohrer invented the STM by measuring quantum tunneling current between an atomic tip and sample surface. \\ \hline
\textbf{32} & \textbf{(a)} & The transmission coefficient decays exponentially with barrier width $a$: $T \approx 16\frac{E}{V_0}(1 - \frac{E}{V_0})e^{-2\kappa a}$. \\ \hline
\textbf{33} & \textbf{(b)} & $[\hat{x}, \hat{p}_x]\psi = x(-i\hbar \psi') - (-i\hbar (x\psi)') = -i\hbar x \psi' + i\hbar \psi + i\hbar x \psi' = i\hbar \psi \implies [\hat{x}, \hat{p}_x] = i\hbar$. \\ \hline
\textbf{34} & \textbf{(a)} & Operators that do not commute do not share a common set of eigenfunctions, establishing a Robertson-Schrödinger uncertainty relation. \\ \hline
\textbf{35} & \textbf{(a)} & Physical observables must be represented by Hermitian operators ($\hat{A}^\dagger = \hat{A}$) to ensure real observable eigenvalues. \\ \hline
\textbf{36} & \textbf{(a)} & Eigenfunctions of Hermitian operators belonging to different eigenvalues are strictly orthogonal. \\ \hline
\textbf{37} & \textbf{(a)} & The quantum expectation value is defined as the sandwich integral $\langle A \rangle = \int \Psi^* \hat{A} \Psi dx$. \\ \hline
\textbf{38} & \textbf{(a)} & The Compton formula is $\Delta\lambda = \lambda_C(1 - \cos\theta)$ where $\lambda_C = \frac{h}{m_e c} \approx 0.02426\text{ \AA} = 2.426\text{ pm}$. \\ \hline
\textbf{39} & \textbf{(a)} & $\lambda_C = \frac{6.626 \times 10^{-34}}{(9.109 \times 10^{-31})(3 \times 10^8)} \approx 2.426 \times 10^{-12}\text{ m} = 0.02426\text{ \AA}$. \\ \hline
\textbf{40} & \textbf{(a)} & At $\theta = 180^\circ$, $(1 - \cos 180^\circ) = 1 - (-1) = 2$, maximizing the wavelength shift to $\Delta\lambda = 2\lambda_C$. \\ \hline
\textbf{41} & \textbf{(a)} & Work function $\Phi = h\nu_0$ is the minimum threshold energy required to overcome the surface barrier potential. \\ \hline
\textbf{42} & \textbf{(a)} & Conservation of photon energy gives $K_{\text{max}} = h\nu - \Phi = e V_0$ where $V_0$ is the stopping potential. \\ \hline
\textbf{43} & \textbf{(a)} & From $V_0 = \left(\frac{h}{e}\right)\nu - \frac{\Phi}{e}$, the slope is universally equal to Planck's constant divided by electron charge ($h/e$). \\ \hline
\textbf{44} & \textbf{(a)} & $\lambda = \frac{1.227}{\sqrt{100}}\text{ nm} = \frac{1.227}{10} = 0.1227\text{ nm} \approx 1.23\text{ \AA}$. \\ \hline
\textbf{45} & \textbf{(a)} & $\lambda = \frac{h}{\sqrt{3 m_n k_B T}} = \frac{6.626 \times 10^{-34}}{\sqrt{3(1.675 \times 10^{-27})(1.38 \times 10^{-23})(300)}} \approx 1.45 \times 10^{-10}\text{ m} = 0.145\text{ nm}$. \\ \hline
\textbf{46} & \textbf{(a)} & $E_1 = \frac{h^2}{8 m_e L^2} = \frac{(6.626 \times 10^{-34})^2}{8(9.109 \times 10^{-31})(10^{-10})^2} \approx 6.024 \times 10^{-18}\text{ J} = \frac{6.024 \times 10^{-18}}{1.602 \times 10^{-19}} \approx 37.6\text{ eV}$. \\ \hline
\textbf{47} & \textbf{(a)} & $E_2 = 2^2 E_1 = 4 \times 37.6\text{ eV} = 150.4\text{ eV}$. \\ \hline
\textbf{48} & \textbf{(a)} & $P = \int_{L/3}^{2L/3}\frac{2}{L}\sin^2(\frac{\pi x}{L})dx = \left[\frac{x}{L} - \frac{\sin(2\pi x/L)}{2\pi}\right]_{L/3}^{2L/3} = \frac{1}{3} + \frac{\sqrt{3}}{2\pi} \approx 0.609$. \\ \hline
\textbf{49} & \textbf{(a)} & $\Delta p_x \ge \frac{\hbar}{2\Delta x} \implies \Delta v_x \ge \frac{\hbar}{2 m_e \Delta x} = \frac{1.0545 \times 10^{-34}}{2(9.109 \times 10^{-31})(10^{-10})} \approx 5.79 \times 10^5\text{ m/s}$. \\ \hline
\textbf{50} & \textbf{(a)} & Ehrenfest's theorem bridges quantum and classical physics, demonstrating $\frac{d\langle\vec{p}\rangle}{dt} = -\langle\nabla V\rangle$. \\ \hline
\end{tabularx}
}

\newpage

\subsection*{Master Answer Key (Questions 51 to 75)}

{\small
\noindent\begin{tabularx}{\textwidth}{|c|c|X|}
\hline
\textbf{Q\#} & \textbf{Ans} & \textbf{Technical Rationale \& Calculation} \\
\hline
\textbf{51} & \textbf{(a)} & $\lambda = \frac{h}{m v} = \frac{6.626 \times 10^{-34}}{(0.01\text{ kg})(10\text{ m/s})} = 6.626 \times 10^{-33}\text{ m}$ (completely unobservable at macroscopic scales). \\ \hline
\textbf{52} & \textbf{(a)} & $p = h/\lambda = m_e c$. Total energy $E = \sqrt{p^2 c^2 + m_e^2 c^4} = \sqrt{2}m_e c^2$. Kinetic energy $K = E - m_e c^2 = (\sqrt{2}-1)m_e c^2 \approx 212\text{ keV}$. \\ \hline
\textbf{53} & \textbf{(a)} & $\Delta E \ge \frac{\hbar}{2\tau} = \frac{1.0545 \times 10^{-34}\text{ J}\cdot\text{s}}{2(10^{-9}\text{ s})} = 5.27 \times 10^{-26}\text{ J} = \frac{5.27 \times 10^{-26}}{1.602 \times 10^{-19}} \approx 3.29 \times 10^{-7}\text{ eV}$. \\ \hline
\textbf{54} & \textbf{(a)} & $\Delta\nu = \frac{\Delta E}{h} = \frac{5.27 \times 10^{-26}}{6.626 \times 10^{-34}} \approx 7.95 \times 10^7\text{ Hz} \approx 80\text{ MHz}$. $\Delta\nu/\nu_0 = \frac{7.95 \times 10^7}{6 \times 10^{14}} \approx 1.33 \times 10^{-7}$. \\ \hline
\textbf{55} & \textbf{(a)} & $\lambda = \frac{h}{\sqrt{2m E_k}} \implies \frac{\lambda_p}{\lambda_e} = \sqrt{\frac{m_e}{m_p}} = \sqrt{\frac{1}{1836}} \approx \frac{1}{42.85}$. \\ \hline
\textbf{56} & \textbf{(a)} & Because $\lambda = h/p$, particles with identical momenta have strictly identical de Broglie wavelengths regardless of mass. \\ \hline
\textbf{57} & \textbf{(a)} & $\Delta E = E_3 - E_1 = 3^2 E_1 - 1^2 E_1 = 8 E_1$. Photon wavelength $\lambda = \frac{hc}{\Delta E} = \frac{hc}{8 E_1}$. \\ \hline
\textbf{58} & \textbf{(a)} & $\langle x^2 \rangle = L^2(\frac{1}{3} - \frac{1}{2\pi^2})$ and $\langle x \rangle^2 = \frac{L^2}{4} \implies \sigma_x^2 = L^2(\frac{1}{12} - \frac{1}{2\pi^2}) \approx 0.0326 L^2$. \\ \hline
\textbf{59} & \textbf{(a)} & $\langle p^2 \rangle = 2m E_1 = \frac{\pi^2\hbar^2}{L^2}$ and $\langle p \rangle = 0 \implies \sigma_p^2 = \frac{\pi^2\hbar^2}{L^2}$. \\ \hline
\textbf{60} & \textbf{(a)} & $\Delta x \Delta p = \sqrt{L^2(\frac{1}{12}-\frac{1}{2\pi^2})}\left(\frac{\pi\hbar}{L}\right) = \hbar\sqrt{\frac{\pi^2}{12}-\frac{1}{2}} \approx 0.568\hbar \ge 0.5\hbar$, validating Heisenberg's principle. \\ \hline
\textbf{61} & \textbf{(a)} & Probability conservation $\frac{\partial P}{\partial t} + \nabla \cdot \vec{J} = 0$ requires $\vec{J} = \frac{\hbar}{2mi}(\Psi^*\nabla\Psi - \Psi\nabla\Psi^*)$. \\ \hline
\textbf{62} & \textbf{(a)} & $\Psi^*\frac{d\Psi}{dx} - \Psi\frac{d\Psi^*}{dx} = A^*(ik A) - A(-ik A^*) = 2ik|A|^2 \implies J = \frac{\hbar}{2mi}(2ik|A|^2) = |A|^2 \frac{\hbar k}{m} = |A|^2 v$. \\ \hline
\textbf{63} & \textbf{(a)} & Since $\psi(x)$ is purely real, $\psi\frac{d\psi}{dx} - \psi\frac{d\psi}{dx} = 0$, confirming zero net probability flow in stationary bound states. \\ \hline
\textbf{64} & \textbf{(a)} & Inside the classically forbidden barrier ($V_0 > E$), the TISE yields real exponential decay $\psi(x) \propto e^{-\kappa x}$ where $\kappa = \sqrt{2m(V_0-E)}/\hbar$. \\ \hline
\textbf{65} & \textbf{(a)} & $\kappa = \frac{\sqrt{2(9.109 \times 10^{-31})(3.0 \times 1.602 \times 10^{-19})}}{1.0545 \times 10^{-34}} = \frac{9.355 \times 10^{-25}}{1.0545 \times 10^{-34}} \approx 8.87 \times 10^9\text{ m}^{-1}$. \\ \hline
\textbf{66} & \textbf{(a)} & $d = \frac{1}{\kappa} = \frac{1}{8.87 \times 10^9\text{ m}^{-1}} \approx 1.127 \times 10^{-10}\text{ m} = 1.13\text{ \AA}$. \\ \hline
\textbf{67} & \textbf{(a)} & Because the transmitted wave is evanescent and carries zero net real current ($J_{\text{trans}} = 0$), all incident probability is reflected ($R = 1.0$). \\ \hline
\textbf{68} & \textbf{(a)} & The quantum harmonic oscillator exhibits equally spaced energy levels $E_n = (n + \frac{1}{2})\hbar\omega$ with zero-point energy $E_0 = \frac{1}{2}\hbar\omega$. \\ \hline
\textbf{69} & \textbf{(a)} & At absolute zero ($n=0$), the ground-state zero-point energy is strictly $E_0 = \frac{1}{2}\hbar\omega$. \\ \hline
\textbf{70} & \textbf{(a)} & Bohr's quantization postulate requires angular momentum to be an integer multiple of Dirac's constant: $L = n\hbar$. \\ \hline
\textbf{71} & \textbf{(a)} & $2\pi r = n\lambda = n(h/p) = n(h/mv) \implies mvr = n(h/2\pi) = n\hbar$, forming constructive standing waves around the circular orbit. \\ \hline
\textbf{72} & \textbf{(a)} & Evaluating with fundamental constants: $a_0 = \frac{4\pi\epsilon_0\hbar^2}{m_e e^2} \approx 0.529177\text{ \AA}$. \\ \hline
\textbf{73} & \textbf{(a)} & Ground-state energy is $E_1 = -13.6\text{ eV}$, so the minimum energy required to remove the electron to infinity is $+13.6\text{ eV}$. \\ \hline
\textbf{74} & \textbf{(a)} & The Rydberg constant for infinite nuclear mass is $R_\infty \approx 1.097373 \times 10^7\text{ m}^{-1}$. \\ \hline
\textbf{75} & \textbf{(a)} & Transitions terminating at ground state $n=1$ ($n \to 1$) form the Lyman series in the ultraviolet ($\lambda = 91\text{--}121\text{ nm}$). \\ \hline
\end{tabularx}
}

\newpage

\subsection*{Master Answer Key (Questions 76 to 100)}

{\small
\noindent\begin{tabularx}{\textwidth}{|c|c|X|}
\hline
\textbf{Q\#} & \textbf{Ans} & \textbf{Technical Rationale \& Calculation} \\
\hline
\textbf{76} & \textbf{(a)} & The Balmer series ($H_\alpha = 656.3\text{ nm}, H_\beta = 486.1\text{ nm}$) spans the visible spectrum. \\ \hline
\textbf{77} & \textbf{(a)} & According to the Born probability postulate, probability of finding eigenvalue $E_1$ is given by the squared modulus $|c_1|^2$. \\ \hline
\textbf{78} & \textbf{(a)} & Symmetric wells yield strictly alternating parity eigenfunctions: $\psi_1(-x) = +\psi_1(x)$ (even), $\psi_2(-x) = -\psi_2(x)$ (odd). \\ \hline
\textbf{79} & \textbf{(a)} & $14 = 1^2 + 2^2 + 3^2$. Six distinct permutations $(1,2,3), (1,3,2), (2,1,3), (2,3,1), (3,1,2), (3,2,1)$ give $6$-fold degeneracy. \\ \hline
\textbf{80} & \textbf{(a)} & $\langle p_x^2 \rangle = 2m E_n = 2m \left(\frac{n^2 \pi^2 \hbar^2}{2m L^2}\right) = \frac{n^2 \pi^2 \hbar^2}{L^2}$. \\ \hline
\textbf{81} & \textbf{(a)} & $|\Psi(x,t)|^2 = \frac{1}{2}|\psi_1|^2 + \frac{1}{2}|\psi_2|^2 + \psi_1\psi_2\cos(\frac{E_2-E_1}{\hbar}t)$, demonstrating quantum beat oscillation. \\ \hline
\textbf{82} & \textbf{(a)} & The Gamow WKB barrier penetration formula is $T \approx \exp(-2\int_{x_1}^{x_2} \kappa(x) dx)$. \\ \hline
\textbf{83} & \textbf{(a)} & Gamow proved that $\alpha$-particles tunnel through the nuclear Coulomb barrier, explaining half-lives spanning $10^{10}\text{ years}$ to $10^{-7}\text{ s}$. \\ \hline
\textbf{84} & \textbf{(a)} & Solving the discontinuity in derivative at $x=0$ yields a single normalized bound state $\psi(x) = \sqrt{\kappa}e^{-\kappa|x|}$ with $E = -\frac{m\alpha^2}{2\hbar^2}$. \\ \hline
\textbf{85} & \textbf{(a)} & Bound states ($E < V_0$) have sinusoidal standing waves inside the well and decaying evanescent tails outside, creating penetrating probability. \\ \hline
\textbf{86} & \textbf{(a)} & In the limit $V_0 \to \infty$, the penetration depth $d = 1/\kappa \to 0$, perfectly recovering the infinite well boundary conditions. \\ \hline
\textbf{87} & \textbf{(a)} & Total wavefunctions of identical fermions must be anti-symmetric under particle exchange: $\Psi(1,2) = -\Psi(2,1)$. \\ \hline
\textbf{88} & \textbf{(a)} & Bosons have symmetric wavefunctions $\Psi(1,2) = +\Psi(2,1)$ and can condense into a single ground state (Bose-Einstein Condensation). \\ \hline
\textbf{89} & \textbf{(a)} & Integrating density of states $g(E) = \frac{V}{2\pi^2}(\frac{2m}{\hbar^2})^{3/2}\sqrt{E}$ up to $E_F$ gives $E_F = \frac{\hbar^2}{2m}(3\pi^2 n)^{2/3}$. \\ \hline
\textbf{90} & \textbf{(a)} & $\langle E \rangle = \frac{1}{N}\int_0^{E_F} E g(E) dE = \frac{\int_0^{E_F} E^{3/2} dE}{\int_0^{E_F} E^{1/2} dE} = \frac{\frac{2}{5}E_F^{5/2}}{\frac{2}{3}E_F^{3/2}} = \frac{3}{5}E_F$. \\ \hline
\textbf{91} & \textbf{(a)} & At $k = \pm n\pi/a$, standing wave formation splits the energy into two distinct levels depending on charge distribution, creating band gaps. \\ \hline
\textbf{92} & \textbf{(a)} & Applying semi-classical acceleration $a = \frac{1}{\hbar}\frac{d v_g}{dt} = \frac{1}{\hbar^2}\frac{d^2 E}{dk^2} F \implies m^* = \frac{\hbar^2}{d^2E/dk^2}$. \\ \hline
\textbf{93} & \textbf{(a)} & Negative curvature implies deceleration under forward electric force; treating missing electrons as positive 'holes' restores standard equations. \\ \hline
\textbf{94} & \textbf{(d)} & Quantum confinement dimensions scale as: 3D ($\sqrt{E}$), 2D (step constant), 1D ($E^{-1/2}$ singularity), 0D (discrete delta functions). \\ \hline
\textbf{95} & \textbf{(a)} & Confinement in all three spatial dimensions collapses energy bands into discrete atomic-like delta function energy levels. \\ \hline
\textbf{96} & \textbf{(a)} & $E_1 = \frac{h^2}{8 m_e^* L^2} = \frac{(6.626 \times 10^{-34})^2}{8(0.067 \times 9.109 \times 10^{-31})(5 \times 10^{-9})^2} \approx 3.60 \times 10^{-21}\text{ J} \approx 22.5\text{ meV}$. \\ \hline
\textbf{97} & \textbf{(a)} & $\nabla^2\psi - \frac{1}{c^2}\frac{\partial^2\psi}{\partial t^2} = \frac{m^2 c^2}{\hbar^2}\psi$ applies to relativistic spin-0 particles. \\ \hline
\textbf{98} & \textbf{(a)} & Paul Dirac's relativistic linear equation $i\hbar\gamma^\mu\partial_\mu\psi = mc\psi$ naturally revealed electron spin and antiparticles. \\ \hline
\textbf{99} & \textbf{(a)} & Electrons passing outside a shielded solenoid acquire a measurable phase shift $\Delta\phi = \frac{q}{\hbar}\oint \vec{A}\cdot d\vec{l} = \frac{q\Phi_B}{\hbar}$. \\ \hline
\textbf{100} & \textbf{(a)} & Entangled particles exhibit non-local quantum correlations violating Bell's inequalities (Einstein-Podolsky-Rosen paradox). \\ \hline
\end{tabularx}
}


\newpage
\section{Section 3: 20 Comprehensive Theory Questions with Analytical Derivations}

\begin{theoryq}{1}
State the de Broglie hypothesis of matter waves. Derive the expression for the de Broglie wavelength of an electron accelerated from rest through a potential difference $V$. Describe the Davisson-Germer experiment and explain how it confirmed matter waves.
\end{theoryq}
\begin{theorysol}
\textbf{1. The de Broglie Hypothesis}:
In 1924, Louis de Broglie postulated that dual wave-particle behavior is a fundamental universal property of all matter. Any moving physical entity of linear momentum $p$ is accompanied by a matter wave of wavelength:
\begin{equation}
\mathbf{\lambda = \frac{h}{p}}
\end{equation}

\textbf{2. Derivation for Electrostatic Acceleration}:
Let an electron of rest mass $m_e$ and charge $e$ be accelerated from rest through a potential difference $V$.
The electrical work done on the electron equals its kinetic energy $E_k$:
\begin{equation}
E_k = e V = \frac{1}{2}m_e v^2 = \frac{p^2}{2m_e} \implies p = \sqrt{2 m_e e V}
\end{equation}
Substituting this into the de Broglie formula:
\begin{equation}
\lambda = \frac{h}{\sqrt{2 m_e e V}}
\end{equation}
Substituting fundamental constants ($h = 6.626 \times 10^{-34}\text{ J}\cdot\text{s}$, $m_e = 9.109 \times 10^{-31}\text{ kg}$, $e = 1.602 \times 10^{-19}\text{ C}$):
\begin{equation}
\mathbf{\lambda = \frac{1.227 \times 10^{-9}\text{ m}}{\sqrt{V}} = \frac{1.227}{\sqrt{V}}\text{ nm} = \frac{12.27}{\sqrt{V}}\text{ \AA}}
\end{equation}

\textbf{3. The Davisson-Germer Experiment}:
\begin{itemize}
    \item \textbf{Setup}: A collimated beam of electrons from a heated tungsten filament was accelerated through variable voltage $V$ onto the cleaved $(111)$ face of a single nickel crystal. Scattered electrons were collected at various scattering angles $\phi$ by a movable Faraday cup detector.
    \item \textbf{Key Observation}: At $V = 54\text{ V}$, an intense, sharp diffraction peak was recorded at scattering angle $\phi = 50^\circ$.
    \item \textbf{Verification via Bragg's Law}:
    The glancing angle with crystal planes (interplanar spacing $d = 0.091\text{ nm}$) was $\theta = 90^\circ - \phi/2 = 90^\circ - 25^\circ = 65^\circ$.
    Bragg's diffraction condition for first order ($n=1$):
    \begin{equation}
    \lambda_{\text{Bragg}} = 2 d \sin\theta = 2(0.091\text{ nm})\sin(65^\circ) = 2(0.091)(0.9063) \approx \mathbf{0.165\text{ nm}}
    \end{equation}
    \item \textbf{Theoretical de Broglie Value}:
    \begin{equation}
    \lambda_{\text{theory}} = \frac{1.227}{\sqrt{54}}\text{ nm} = \frac{1.227}{7.348} \approx \mathbf{0.167\text{ nm}}
    \end{equation}
    The experimental Bragg wavelength ($0.165\text{ nm}$) matched de Broglie's prediction ($0.167\text{ nm}$) to within $1\%$, providing direct experimental proof of matter waves.
\end{itemize}
\end{theorysol}
\vspace{4mm}

\begin{theoryq}{2}
Define Phase Velocity ($v_p$) and Group Velocity ($v_g$). Derive the general relationship between group velocity and particle velocity ($v_g = v_{\text{particle}}$) for both non-relativistic and relativistic matter waves. Show that $v_p \cdot v_g = c^2$.
\end{theoryq}
\begin{theorysol}
\textbf{1. Definitions}:
\begin{itemize}
    \item \textbf{Phase Velocity ($v_p$)}: The speed at which individual monochromatic wave crests and troughs propagate through space:
    \begin{equation}
    \mathbf{v_p = \frac{\omega}{k}}
    \end{equation}
    \item \textbf{Group Velocity ($v_g$)}: The speed at which the overall envelope modulating a localized wave packet propagates through space (the velocity of energy and information transport):
    \begin{equation}
    \mathbf{v_g = \frac{d\omega}{dk}}
    \end{equation}
\end{itemize}

\textbf{2. Proof of $v_g = v_{\text{particle}}$ (Non-Relativistic)}:
For a free particle of mass $m$ and velocity $v$:
Total energy $E = \hbar\omega = \frac{p^2}{2m} = \frac{\hbar^2 k^2}{2m} \implies \omega(k) = \frac{\hbar k^2}{2m}$.
Differentiating with respect to $k$:
\begin{equation}
\mathbf{v_g = \frac{d\omega}{dk} = \frac{\hbar(2k)}{2m} = \frac{\hbar k}{m} = \frac{p}{m} = v_{\text{particle}}}
\end{equation}
Phase velocity is $v_p = \frac{\omega}{k} = \frac{\hbar k}{2m} = \frac{v_{\text{particle}}}{2}$.

\textbf{3. Proof of $v_g = v_{\text{particle}}$ \& $v_p \cdot v_g = c^2$ (Relativistic)}:
According to special relativity:
\begin{equation}
E^2 = p^2 c^2 + m_0^2 c^4 \implies \hbar^2 \omega^2 = \hbar^2 k^2 c^2 + m_0^2 c^4
\end{equation}
Differentiating both sides with respect to $k$:
\begin{equation}
2\hbar^2 \omega \frac{d\omega}{dk} = 2\hbar^2 k c^2 \implies \frac{d\omega}{dk} = c^2 \left(\frac{k}{\omega}\right) = \frac{c^2}{v_p}
\end{equation}
Rearranging:
\begin{equation}
\mathbf{v_p \cdot v_g = c^2}
\end{equation}
Since $E = \hbar\omega = \gamma m_0 c^2$ and $p = \hbar k = \gamma m_0 v$:
\begin{align}
v_p &= \frac{\omega}{k} = \frac{\hbar\omega}{\hbar k} = \frac{\gamma m_0 c^2}{\gamma m_0 v} = \frac{c^2}{v} \\
v_g &= \frac{c^2}{v_p} = \frac{c^2}{c^2/v} = \mathbf{v_{\text{particle}}}
\end{align}
\textit{Conclusion}: The group velocity of the quantum wave packet moves at the physical particle speed $v$, while the phase velocity $v_p = c^2/v > c$ exceeds light speed without violating causality.
\end{theorysol}
\vspace{4mm}

\begin{theoryq}{3}
State Heisenberg's Uncertainty Principle for position and momentum. Use it to prove analytically that electrons cannot reside inside an atomic nucleus. Estimate the ground-state energy of a hydrogen atom using the uncertainty principle.
\end{theoryq}
\begin{theorysol}
\textbf{1. Statement}:
Heisenberg's Uncertainty Principle states that it is physically impossible to measure simultaneously the exact position $x$ and linear momentum $p_x$ of a particle with unlimited precision:
\begin{equation}
\mathbf{\Delta x \cdot \Delta p_x \ge \frac{\hbar}{2}}
\end{equation}

\textbf{2. Proof of Non-Existence of Electrons Inside the Nucleus}:
Typical nuclear diameter is $D \sim 10^{-14}\text{ m}$. If an electron exists inside the nucleus, its position uncertainty is at most $\Delta x \approx 10^{-14}\text{ m}$.
By Heisenberg's principle:
\begin{equation}
\Delta p_x \ge \frac{\hbar}{2\Delta x} \approx \frac{1.0545 \times 10^{-34}\text{ J}\cdot\text{s}}{2(10^{-14}\text{ m})} \approx 5.27 \times 10^{-21}\text{ kg}\cdot\text{m/s}
\end{equation}
Because momentum must be at least comparable to uncertainty ($p \ge \Delta p_x$):
Calculating relativistic kinetic energy:
\begin{align*}
E &\approx p c = (5.27 \times 10^{-21}\text{ kg}\cdot\text{m/s})(3 \times 10^8\text{ m/s}) \approx 1.58 \times 10^{-12}\text{ J} \\
E &= \frac{1.58 \times 10^{-12}\text{ J}}{1.602 \times 10^{-19}\text{ J/eV}} \approx 9.87 \times 10^6\text{ eV} \approx \mathbf{10\text{--}20\text{ MeV}}
\end{align*}
However, beta-decay electrons emitted from radioactive nuclei possess energies of only $2\text{--}4\text{ MeV}$. If electrons were permanently bound inside the nucleus, their energies would have to exceed $20\text{ MeV}$. This contradiction proves that electrons cannot exist inside the nucleus.

\textbf{3. Ground-State Energy of Hydrogen Atom}:
Let the electron be confined within distance $r$ from the proton ($\Delta x \approx r$).
Then momentum uncertainty is $p \approx \Delta p \approx \frac{\hbar}{r}$.
Total energy of the hydrogen atom is:
\begin{equation}
E(r) = E_k + V = \frac{p^2}{2m} - \frac{e^2}{4\pi\epsilon_0 r} = \frac{\hbar^2}{2m r^2} - \frac{e^2}{4\pi\epsilon_0 r}
\end{equation}
To find the stable ground state, minimize $E(r)$ with respect to $r$:
\begin{equation}
\frac{dE}{dr} = -\frac{\hbar^2}{m r^3} + \frac{e^2}{4\pi\epsilon_0 r^2} = 0 \implies \mathbf{r_0 = \frac{4\pi\epsilon_0 \hbar^2}{m e^2} = a_0 \approx 0.529\text{ \AA}}
\end{equation}
Substituting $r_0$ back into $E(r)$ gives the exact ground state energy:
\begin{equation}
\mathbf{E_1 = -\frac{m e^4}{32\pi^2\epsilon_0^2\hbar^2} = -13.6\text{ eV}}
\end{equation}
\end{theorysol}
\vspace{4mm}

\begin{theoryq}{4}
Explain the physical significance of the quantum mechanical wave function $\Psi(x,t)$. Detail Max Born's probability interpretation, the normalization condition, and the four mathematical boundary conditions for physically acceptable wave functions.
\end{theoryq}
\begin{theorysol}
\textbf{1. Physical Interpretation of $\Psi(x,t)$}:
In quantum mechanics, the state of a physical system is completely specified by a complex wave function $\Psi(x,t)$.
While $\Psi$ itself has no direct classical counterpart and is not directly measurable, its squared magnitude carries direct physical meaning.

\textbf{2. Max Born's Statistical Postulate}:
The probability of finding a particle in an infinitesimal spatial interval $dx$ at time $t$ is proportional to the square of the modulus of the wave function:
\begin{equation}
\mathbf{dP(x,t) = |\Psi(x,t)|^2 dx = \Psi^*(x,t)\Psi(x,t) dx}
\end{equation}
where $P(x,t) = |\Psi(x,t)|^2$ is the \textbf{Probability Density}.

\textbf{3. Normalization Condition}:
Because the particle must exist somewhere along the 1D space with certainty ($100\%$ probability):
\begin{equation}
\mathbf{\int_{-\infty}^{+\infty} |\Psi(x,t)|^2 dx = 1}
\end{equation}

\textbf{4. Four Mathematical Boundary Conditions for Acceptable $\Psi(x)$}:
To ensure physical realism and conservation of probability, $\psi(x)$ must satisfy:
\begin{enumerate}
    \item \textbf{Single-Valued}: $\psi(x)$ must have only one unique value at any given spatial coordinate $x$ so that probability density $|\psi(x)|^2$ is unambiguously defined.
    \item \textbf{Continuous Everywhere}: $\psi(x)$ must be spatially continuous across all boundaries ($\lim_{\epsilon\to 0}\psi(x_0+\epsilon) = \lim_{\epsilon\to 0}\psi(x_0-\epsilon)$).
    \item \textbf{Continuous First Derivative}: $\frac{d\psi}{dx}$ must be continuous everywhere for all finite potential potentials $V(x)$, ensuring momentum is uniquely defined.
    \item \textbf{Square-Integrable (Finite Everywhere)}: $\psi(x)$ must vanish as $x \to \pm\infty$ ($\lim_{x\to\pm\infty}\psi(x) = 0$) so that the total probability integral remains finite and normalizable.
\end{enumerate}

\end{theorysol}
\vspace{4mm}

\begin{theoryq}{5}
Derive the 1D Time-Independent Schrödinger Equation (TISE) from the Time-Dependent Schrödinger Equation (TDSE) by the method of separation of variables. Define stationary states.
\end{theoryq}
\begin{theorysol}
\textbf{1. Time-Dependent Schrödinger Equation (TDSE)}:
For a particle of mass $m$ in a time-independent potential $V(x)$:
\begin{equation}
i\hbar \frac{\partial \Psi(x,t)}{\partial t} = -\frac{\hbar^2}{2m}\frac{\partial^2 \Psi(x,t)}{\partial x^2} + V(x)\Psi(x,t)
\end{equation}

\textbf{2. Separation of Variables}:
Assume a separable product solution:
\begin{equation}
\Psi(x,t) = \psi(x) \phi(t)
\end{equation}
Compute partial derivatives:
\begin{equation}
\frac{\partial \Psi}{\partial t} = \psi(x)\frac{d\phi(t)}{dt}, \qquad \frac{\partial^2 \Psi}{\partial x^2} = \phi(t)\frac{d^2\psi(x)}{dx^2}
\end{equation}
Substitute these into the TDSE:
\begin{equation}
i\hbar \psi(x)\frac{d\phi(t)}{dt} = -\frac{\hbar^2}{2m}\phi(t)\frac{d^2\psi(x)}{dx^2} + V(x)\psi(x)\phi(t)
\end{equation}
Divide the entire equation by $\Psi(x,t) = \psi(x)\phi(t)$:
\begin{equation}
i\hbar \frac{1}{\phi(t)}\frac{d\phi(t)}{dt} = \frac{1}{\psi(x)}\left[ -\frac{\hbar^2}{2m}\frac{d^2\psi(x)}{dx^2} + V(x)\psi(x) \right]
\end{equation}
The left side depends solely on time $t$, while the right side depends solely on position $x$. For this to hold for all $x$ and $t$, both sides must equal a common separation constant $E$ (Total Energy):

\textbf{3. Spatial Equation (TISE)}:
\begin{equation}
\mathbf{-\frac{\hbar^2}{2m}\frac{d^2\psi(x)}{dx^2} + V(x)\psi(x) = E\psi(x) \iff \hat{H}\psi = E\psi}
\end{equation}

\textbf{4. Temporal Equation \& Stationary States}:
\begin{equation}
i\hbar \frac{d\phi(t)}{dt} = E\phi(t) \implies \frac{d\phi}{\phi} = -\frac{i E}{\hbar}dt \implies \mathbf{\phi(t) = e^{-i E t/\hbar}}
\end{equation}
The full wavefunction is $\Psi(x,t) = \psi(x)e^{-iEt/\hbar}$.
\textit{Stationary State}: The probability density is strictly static in time:
\begin{equation}
P(x,t) = |\Psi(x,t)|^2 = |\psi(x)|^2 |e^{-iEt/\hbar}|^2 = |\psi(x)|^2
\end{equation}
\end{theorysol}
\vspace{4mm}

\begin{theoryq}{6}
Solve the Time-Independent Schrödinger Equation for a particle of mass $m$ trapped in a 1D Infinite Potential Well of width $L$ ($V(x) = 0$ for $0 < x < L$ and $V(x) = \infty$ elsewhere). Derive the quantized energy eigenvalues $E_n$ and normalized eigenfunctions $\psi_n(x)$.
\end{theoryq}
\begin{theorysol}
\textbf{1. Model Formulation \& Boundary Conditions}:
Potential profile:
\begin{equation}
V(x) = \begin{cases} 0 & 0 < x < L \\ \infty & x \le 0 \text{ or } x \ge L \end{cases}
\end{equation}
Since $V = \infty$ outside the well, the particle cannot penetrate outside: $\psi(x) = 0$ for $x \le 0$ and $x \ge L$.
Continuity of $\psi$ requires boundary conditions:
\begin{equation}
\psi(0) = 0 \quad \text{and} \quad \psi(L) = 0
\end{equation}

\textbf{2. General Solution Inside the Well ($0 < x < L$)}:
Since $V(x) = 0$, the TISE is:
\begin{equation}
-\frac{\hbar^2}{2m}\frac{d^2\psi}{dx^2} = E\psi \implies \frac{d^2\psi}{dx^2} + k^2\psi = 0
\end{equation}
where $k = \frac{\sqrt{2mE}}{\hbar}$.
The general solution is:
\begin{equation}
\psi(x) = A\sin(kx) + B\cos(kx)
\end{equation}

\textbf{3. Applying Boundary Conditions}:
\begin{enumerate}
    \item At $x = 0$: $\psi(0) = A\sin(0) + B\cos(0) = B = 0 \implies B = 0$.
    \item Thus $\psi(x) = A\sin(kx)$.
    \item At $x = L$: $\psi(L) = A\sin(kL) = 0$.
    Since $A \ne 0$ (otherwise $\psi=0$ identically), we must have:
    \begin{equation}
    k L = n\pi \implies \mathbf{k_n = \frac{n\pi}{L} \quad (n = 1, 2, 3,\dots)}
    \end{equation}
\end{enumerate}

\textbf{4. Quantized Energy Eigenvalues $E_n$}:
\begin{equation}
E_n = \frac{\hbar^2 k_n^2}{2m} = \frac{\hbar^2}{2m}\left(\frac{n\pi}{L}\right)^2 = \mathbf{\frac{n^2 \pi^2 \hbar^2}{2m L^2} = \frac{n^2 h^2}{8m L^2}} \quad (n = 1, 2, 3,\dots)
\end{equation}

\textbf{5. Normalization of Eigenfunctions}:
\begin{align*}
\int_0^L |\psi_n(x)|^2 dx = 1 &\implies A^2 \int_0^L \sin^2\left(\frac{n\pi x}{L}\right)dx = 1 \\
A^2 \int_0^L \frac{1 - \cos(2n\pi x/L)}{2}dx &= A^2 \left(\frac{L}{2}\right) = 1 \implies \mathbf{A = \sqrt{\frac{2}{L}}}
\end{align*}
The complete \textbf{Normalized Wavefunctions} are:
\begin{equation}
\mathbf{\psi_n(x) = \sqrt{\frac{2}{L}}\sin\left(\frac{n\pi x}{L}\right)}
\end{equation}
\end{theorysol}
\vspace{4mm}

\begin{theoryq}{7}
For a particle in the ground state ($n=1$) of a 1D box of width $L$, calculate the expectation values $\langle x \rangle, \langle x^2 \rangle, \langle p_x \rangle,$ and $\langle p_x^2 \rangle$. Verify that $\Delta x \cdot \Delta p_x \ge \hbar/2$.
\end{theoryq}
\begin{theorysol}
\textbf{1. Wavefunction}:
$\psi_1(x) = \sqrt{\frac{2}{L}}\sin\left(\frac{\pi x}{L}\right)$ for $0 \le x \le L$.

\textbf{2. Position Expectation Values}:
\begin{itemize}
    \item $\langle x \rangle$:
    \begin{equation}
    \langle x \rangle = \int_0^L \psi_1^* x \psi_1 dx = \frac{2}{L}\int_0^L x \sin^2\left(\frac{\pi x}{L}\right)dx = \mathbf{\frac{L}{2}}
    \end{equation}
    \item $\langle x^2 \rangle$:
    Using identity $\sin^2\theta = \frac{1-\cos 2\theta}{2}$ and integrating by parts:
    \begin{equation}
    \langle x^2 \rangle = \frac{2}{L}\int_0^L x^2 \left[\frac{1 - \cos(2\pi x/L)}{2}\right]dx = \mathbf{L^2 \left( \frac{1}{3} - \frac{1}{2\pi^2} \right)} \approx 0.2827 L^2
    \end{equation}
    \item Position uncertainty $\Delta x$:
    \begin{equation}
    (\Delta x)^2 = \langle x^2 \rangle - \langle x \rangle^2 = L^2\left(\frac{1}{3} - \frac{1}{2\pi^2}\right) - \frac{L^2}{4} = L^2\left(\frac{1}{12} - \frac{1}{2\pi^2}\right) \approx \mathbf{0.03267 L^2} \implies \Delta x \approx 0.1807 L
    \end{equation}
\end{itemize}

\textbf{3. Momentum Expectation Values}:
\begin{itemize}
    \item $\langle p_x \rangle$:
    \begin{equation}
    \langle p_x \rangle = \int_0^L \psi_1^* \left(-i\hbar\frac{d}{dx}\right)\psi_1 dx = -\frac{2i\hbar\pi}{L^2}\int_0^L \sin\left(\frac{\pi x}{L}\right)\cos\left(\frac{\pi x}{L}\right)dx = \mathbf{0}
    \end{equation}
    \item $\langle p_x^2 \rangle$:
    \begin{equation}
    \langle p_x^2 \rangle = \int_0^L \psi_1^* \left(-\hbar^2\frac{d^2}{dx^2}\right)\psi_1 dx = \frac{\pi^2\hbar^2}{L^2}\int_0^L \psi_1^2 dx = \mathbf{\frac{\pi^2\hbar^2}{L^2}}
    \end{equation}
    \item Momentum uncertainty $\Delta p_x$:
    \begin{equation}
    \Delta p_x = \sqrt{\langle p_x^2 \rangle - \langle p_x \rangle^2} = \mathbf{\frac{\pi\hbar}{L}}
    \end{equation}
\end{itemize}

\textbf{4. Uncertainty Product Verification}:
\begin{equation}
\mathbf{\Delta x \cdot \Delta p_x = (0.1807 L)\left(\frac{\pi\hbar}{L}\right) = 0.1807 \pi \hbar \approx 0.568 \hbar}
\end{equation}
Since $0.568\hbar > 0.500\hbar$ ($\hbar/2$), Heisenberg's Uncertainty Principle is rigorously satisfied.
\end{theorysol}
\vspace{4mm}

\begin{theoryq}{8}
Extend the Schrödinger equation to a 3D rectangular box ($a \times b \times c$) and cubical box ($L \times L \times L$). Derive the energy eigenvalues and explain the physical concept of energy degeneracy with specific examples.
\end{theoryq}
\begin{theorysol}
\textbf{1. 3D Potential Well \& Separation of Variables}:
Potential $V(x,y,z) = 0$ inside $0 < x < a, 0 < y < b, 0 < z < c$, and $\infty$ elsewhere.
The 3D TISE is:
\begin{equation}
-\frac{\hbar^2}{2m}\left(\frac{\partial^2}{\partial x^2} + \frac{\partial^2}{\partial y^2} + \frac{\partial^2}{\partial z^2}\right)\psi(x,y,z) = E\psi(x,y,z)
\end{equation}
Substituting $\psi(x,y,z) = X(x)Y(y)Z(z)$ and $E = E_x + E_y + E_z$:
Applying boundary conditions at each face yields:
\begin{equation}
\mathbf{\psi_{n_x,n_y,n_z}(x,y,z) = \sqrt{\frac{8}{a b c}}\sin\left(\frac{n_x\pi x}{a}\right)\sin\left(\frac{n_y\pi y}{b}\right)\sin\left(\frac{n_z\pi z}{c}\right)}
\end{equation}
where $n_x, n_y, n_z \in \{1, 2, 3,\dots\}$.

\textbf{2. Energy Eigenvalues}:
\begin{equation}
\mathbf{E_{n_x,n_y,n_z} = \frac{h^2}{8m}\left( \frac{n_x^2}{a^2} + \frac{n_y^2}{b^2} + \frac{n_z^2}{c^2} \right)}
\end{equation}

\textbf{3. Cubical Box ($a = b = c = L$)}:
\begin{equation}
\mathbf{E = \frac{h^2}{8mL^2}(n_x^2 + n_y^2 + n_z^2) = E_1(n_x^2 + n_y^2 + n_z^2)}
\end{equation}
where $E_1 = \frac{h^2}{8mL^2}$.

\textbf{4. Concept of Energy Degeneracy}:
Degeneracy is the condition where two or more linearly independent quantum eigenfunctions share the exact same energy eigenvalue:
\begin{itemize}
    \item \textbf{Ground State}: $(1,1,1) \implies E = 3E_1$ (\textbf{Non-degenerate}, $g=1$).
    \item \textbf{First Excited State}: $n_x^2+n_y^2+n_z^2 = 1^2+1^2+2^2 = 6 \implies E = 6E_1$.
    Three distinct states: $(2,1,1), (1,2,1), (1,1,2)$ (\textbf{3-fold degenerate}, $g=3$).
    \item \textbf{Second Excited State}: $n_x^2+n_y^2+n_z^2 = 1^2+2^2+2^2 = 9 \implies E = 9E_1$.
    Three distinct states: $(1,2,2), (2,1,2), (2,2,1)$ (\textbf{3-fold degenerate}, $g=3$).
    \item \textbf{Third Excited State}: $n_x^2+n_y^2+n_z^2 = 1^2+1^2+3^2 = 11 \implies E = 11E_1$.
    Three states: $(3,1,1), (1,3,1), (1,1,3)$ ($g=3$).
    \item \textbf{State $E = 14E_1$}: $1^2+2^2+3^2 = 14$.
    Six states: $(1,2,3), (1,3,2), (2,1,3), (2,3,1), (3,1,2), (3,2,1)$ (\textbf{6-fold degenerate}, $g=6$).
\end{itemize}
\end{theorysol}
\vspace{4mm}

\begin{theoryq}{9}
Explain the physical principle of Quantum Mechanical Tunneling. Set up the Schrödinger equation for a rectangular potential barrier of height $V_0$ and width $a$ for $E < V_0$. State the transmission coefficient formula and explain its application in the Scanning Tunneling Microscope (STM).
\end{theoryq}
\begin{theorysol}
\textbf{1. Concept of Quantum Tunneling}:
Classically, a particle of total energy $E$ encountering a potential barrier $V_0 > E$ cannot penetrate into the barrier and must undergo $100\%$ reflection.
In quantum mechanics, the wave function does not drop abruptly to zero at the barrier edge but decays exponentially: $\psi(x) \propto e^{-\kappa x}$. If the barrier is thin, $\psi(x)$ remains non-zero at the exit boundary, emerging as a transmitted travelling wave with finite probability $T > 0$.

\textbf{2. Boundary Value Formulation}:
Let a barrier of width $a$ extend from $x = 0$ to $x = a$:
\begin{align*}
\text{Region I } (x < 0, V=0): &\quad \psi_{\text{I}}(x) = e^{i k_1 x} + R e^{-i k_1 x}, \quad k_1 = \frac{\sqrt{2mE}}{\hbar} \\
\text{Region II } (0 < x < a, V=V_0): &\quad \psi_{\text{II}}(x) = A e^{-\kappa x} + B e^{+\kappa x}, \quad \kappa = \frac{\sqrt{2m(V_0 - E)}}{\hbar} \\
\text{Region III } (x > a, V=0): &\quad \psi_{\text{III}}(x) = T e^{i k_1 x}
\end{align*}

\textbf{3. Transmission Coefficient $T$}:
Matching $\psi$ and $\frac{d\psi}{dx}$ at $x = 0$ and $x = a$ yields:
\begin{equation}
\mathbf{T = \left[ 1 + \frac{V_0^2 \sinh^2(\kappa a)}{4E(V_0 - E)} \right]^{-1} \approx 16 \frac{E}{V_0}\left(1 - \frac{E}{V_0}\right)e^{-2\kappa a} \quad (\text{for } \kappa a \gg 1)}
\end{equation}

\textbf{4. Application: Scanning Tunneling Microscope (STM)}:
\begin{itemize}
    \item An atomically sharp metallic tip (Tungsten or Pt-Ir) is brought within sub-nanometer distance ($s \approx 0.5\text{ nm}$) of a conducting sample.
    \item Applying a small bias voltage $V$ induces a tunneling current $I_{\text{tunnel}} \propto e^{-2\kappa s}$.
    \item Because $I_{\text{tunnel}}$ changes by an order of magnitude ($10\times$) for every $0.1\text{ nm}$ change in tip-sample gap $s$, piezo-electric scanners can map individual surface atoms with sub-angstrom vertical resolution.
\end{itemize}
\end{theorysol}
\vspace{4mm}

\begin{theoryq}{10}
An electron is confined within a 1D quantum well of width $L = 0.50\text{ nm}$. Calculate: (a) the first three energy levels $E_1, E_2, E_3$ in $\text{eV}$, (b) the wavelength of light emitted when the electron transitions from $n=3$ to $n=1$, (c) the probability of finding the electron in the range $0 \le x \le L/4$ in the ground state.
\end{theoryq}
\begin{theorysol}
\textbf{1. Energy Levels $E_1, E_2, E_3$}:
\begin{align*}
E_1 &= \frac{h^2}{8 m_e L^2} = \frac{(6.626 \times 10^{-34}\text{ J}\cdot\text{s})^2}{8(9.109 \times 10^{-31}\text{ kg})(0.50 \times 10^{-9}\text{ m})^2} \\
&= \frac{4.3904 \times 10^{-67}}{1.8218 \times 10^{-48}} \approx 2.4099 \times 10^{-19}\text{ J} = \frac{2.4099 \times 10^{-19}}{1.602 \times 10^{-19}} \approx \mathbf{1.504\text{ eV}} \\
\mathbf{E_2} &= 2^2 E_1 = 4(1.504\text{ eV}) = \mathbf{6.016\text{ eV}} \\
\mathbf{E_3} &= 3^2 E_1 = 9(1.504\text{ eV}) = \mathbf{13.536\text{ eV}}
\end{align*}

\textbf{2. Transition Wavelength ($n=3 \to n=1$)}:
\begin{align*}
\Delta E &= E_3 - E_1 = 13.536 - 1.504 = 12.032\text{ eV} = 1.9275 \times 10^{-18}\text{ J} \\
\mathbf{\lambda} &= \frac{h c}{\Delta E} = \frac{(6.626 \times 10^{-34})(3 \times 10^8)}{1.9275 \times 10^{-18}} \approx 1.031 \times 10^{-7}\text{ m} = \mathbf{103.1\text{ nm}} \quad (\text{Far UV})
\end{align*}

\textbf{3. Probability in $0 \le x \le L/4$ (Ground State)}:
\begin{align*}
P &= \int_0^{L/4} \frac{2}{L}\sin^2\left(\frac{\pi x}{L}\right)dx = \frac{2}{L}\int_0^{L/4} \left(\frac{1 - \cos(2\pi x/L)}{2}\right)dx \\
&= \frac{1}{L}\left[ x - \frac{L}{2\pi}\sin\left(\frac{2\pi x}{L}\right) \right]_0^{L/4} = \frac{1}{L}\left[ \frac{L}{4} - \frac{L}{2\pi}\sin\left(\frac{\pi}{2}\right) \right] \\
&= \frac{1}{4} - \frac{1}{2\pi} = 0.2500 - 0.1592 = \mathbf{0.0908 = 9.08\%}
\end{align*}
\end{theorysol}
\vspace{4mm}

\begin{theoryq}{11}
Describe the Davisson-Germer Experiment. Explain the experimental setup, working method, observation of diffraction peaks, and show how the calculated de Broglie wavelength matches Bragg's law.
\end{theoryq}
\begin{theorysol}
\textbf{1. Experimental Setup and Working Principle}:
\begin{itemize}
\item An electron gun produces a collimated beam of electrons accelerated by an adjustable potential $V$ ($40\text{--}68\text{ V}$).
\item The beam strikes the target surface of a single crystal Nickel (Ni) target.
\item Scattered electrons are collected by a movable electrostatic Faraday collector connected to a sensitive galvanometer at various scattering angles $\phi$.
\end{itemize}

\textbf{2. Experimental Observations}:
\begin{itemize}
\item At an accelerating voltage $V = 54\text{ V}$, a prominent, sharp intensity peak occurs at a scattering angle $\phi = 50^\circ$.
\item The glancing angle (Bragg angle) $\theta$ with respect to the atomic crystal planes is:
\begin{equation}
\theta = 90^\circ - \frac{\phi}{2} = 90^\circ - 25^\circ = 65^\circ
\end{equation}
\end{itemize}

\textbf{3. Verification using Bragg's Law}:
For Nickel, interplanar spacing from X-ray diffraction is $d = 0.091\text{ nm} = 0.91\text{ \AA}$.
Using first-order Bragg's diffraction condition ($n=1$):
\begin{equation}
\lambda_{Bragg} = 2d \sin\theta = 2(0.091\text{ nm}) \sin(65^\circ) = 2(0.091)(0.9063) = \mathbf{0.165\text{ nm}}
\end{equation}

\textbf{4. Theoretical de Broglie Wavelength}:
Using de Broglie's formula for an electron accelerated through $V = 54\text{ V}$:
\begin{equation}
\lambda_{dB} = \frac{h}{\sqrt{2 m_e e V}} = \frac{1.227}{\sqrt{54}}\text{ nm} = \frac{1.227}{7.348}\text{ nm} = \mathbf{0.167\text{ nm}}
\end{equation}
The near-perfect agreement between $\lambda_{Bragg}$ ($0.165\text{ nm}$) and $\lambda_{dB}$ ($0.167\text{ nm}$) provided the first conclusive experimental proof of de Broglie's wave hypothesis.
\end{theorysol}
\vspace{4mm}

\begin{theoryq}{12}
Define Hermitian Operators in quantum mechanics. Prove that the eigenvalues of Hermitian operators are strictly real and that eigenfunctions corresponding to distinct eigenvalues are mutually orthogonal.
\end{theoryq}
\begin{theorysol}
\textbf{1. Definition of Hermitian Operator}:
An operator $\hat{A}$ is Hermitian if for any well-behaved wavefunctions $\psi$ and $\phi$:
\begin{equation}
\int \psi^* (\hat{A} \phi)\, dx = \int (\hat{A}\psi)^* \phi\, dx = \langle \psi | \hat{A} | \phi \rangle = \langle \hat{A}\psi | \phi \rangle
\end{equation}
Physical observables (position, momentum, energy) are represented by Hermitian operators because physical measurements must yield real values.

\textbf{2. Proof that Eigenvalues are Strictly Real}:
Let $\hat{A}\psi_n = a_n \psi_n$.
Taking the inner product with $\psi_n$:
\begin{equation}
\langle \psi_n | \hat{A} | \psi_n \rangle = \int \psi_n^* (\hat{A} \psi_n)\, dx = a_n \int |\psi_n|^2 dx = a_n \langle \psi_n | \psi_n \rangle
\end{equation}
Using the Hermiticity property:
\begin{equation}
\langle \psi_n | \hat{A} | \psi_n \rangle = \langle \hat{A}\psi_n | \psi_n \rangle = \int (\hat{A}\psi_n)^* \psi_n\, dx = a_n^* \int |\psi_n|^2 dx = a_n^* \langle \psi_n | \psi_n \rangle
\end{equation}
Subtracting the two relations:
\begin{equation}
(a_n - a_n^*) \langle \psi_n | \psi_n \rangle = 0
\end{equation}
Since $\langle \psi_n | \psi_n \rangle > 0$ for non-trivial states, $a_n = a_n^*$, proving all eigenvalues $a_n$ are purely real.

\textbf{3. Proof of Orthogonality of Eigenfunctions}:
Let $\hat{A}\psi_m = a_m \psi_m$ and $\hat{A}\psi_n = a_n \psi_n$ with $a_m \ne a_n$.
\begin{equation}
\langle \psi_m | \hat{A} | \psi_n \rangle = a_n \langle \psi_m | \psi_n \rangle
\end{equation}
\begin{equation}
\langle \hat{A}\psi_m | \psi_n \rangle = a_m^* \langle \psi_m | \psi_n \rangle = a_m \langle \psi_m | \psi_n \rangle
\end{equation}
Subtracting:
\begin{equation}
(a_n - a_m) \langle \psi_m | \psi_n \rangle = 0
\end{equation}
Since $a_m \ne a_n$, we must have $\langle \psi_m | \psi_n \rangle = \int \psi_m^* \psi_n dx = 0$, proving distinct eigenfunctions are orthogonal.
\end{theorysol}
\vspace{4mm}

\begin{theoryq}{13}
Solve the Schrödinger equation for a 1D Quantum Harmonic Oscillator. State the quantized energy eigenvalues, explain zero-point energy $E_0 = \frac{1}{2}\hbar\omega$, and compare quantum vs classical probability distributions.
\end{theoryq}
\begin{theorysol}
\textbf{1. Formulation of the Oscillator Equation}:
For a particle of mass $m$ experiencing restoring force $F = -k x$, potential energy is $V(x) = \frac{1}{2}m\omega^2 x^2$.
The Time-Independent Schrödinger Equation is:
\begin{equation}
-\frac{\hbar^2}{2m}\frac{d^2\psi}{dx^2} + \frac{1}{2}m\omega^2 x^2 \psi = E\psi
\end{equation}

\textbf{2. Dimensionless Form and Hermite Polynomials}:
Let $y = \sqrt{\frac{m\omega}{\hbar}} x$ and $\epsilon = \frac{2E}{\hbar\omega}$.
The asymptotic solution is $\psi(y) = H(y) e^{-y^2/2}$, where $H(y)$ satisfies Hermite's differential equation.
For the wavefunction to remain square-integrable as $y \to \pm\infty$, the power series must terminate into a polynomial of order $n$, yielding:
\begin{equation}
\epsilon = 2n + 1 \implies \mathbf{E_n = \left(n + \frac{1}{2}\right)\hbar\omega},\quad n = 0, 1, 2, \dots
\end{equation}
Normalized eigenfunctions are:
\begin{equation}
\psi_n(x) = \frac{1}{\sqrt{2^n n!}}\left(\frac{m\omega}{\pi\hbar}\right)^{1/4} H_n\left(\sqrt{\frac{m\omega}{\hbar}}x\right) e^{-\frac{m\omega x^2}{2\hbar}}
\end{equation}

\textbf{3. Zero-Point Energy}:
Even at absolute zero ($n=0$), ground state energy is $E_0 = \frac{1}{2}\hbar\omega > 0$.
This is a direct consequence of Heisenberg's Uncertainty Principle: if $E_0 = 0$, the particle would sit motionless at $x=0$ ($\Delta x = 0, \Delta p = 0$), violating $\Delta x \Delta p \ge \hbar/2$.

\textbf{4. Classical vs Quantum Probability}:
\begin{itemize}
\item \textit{Classical:} Maximum probability density occurs at classical turning points ($x = \pm A$) where velocity $v \to 0$.
\item \textit{Quantum Ground State ($n=0$):} Maximum probability density $|\psi_0(x)|^2$ occurs at the center ($x=0$), with exponential tunneling into classically forbidden regions $|x| > A$.
\item \textit{Correspondence Principle:} For large quantum numbers $n \gg 1$, quantum probability peaks match classical turning points.
\end{itemize}

\end{theorysol}
\vspace{4mm}

\begin{theoryq}{14}
Analyze a 1D Finite Potential Well of depth $V_0$ and width $2a$. Explain the existence of bound states ($E < V_0$), boundary conditions, and show why the number of bound states is finite.
\end{theoryq}
\begin{theorysol}
\textbf{1. Potential Profile}:
\begin{equation}
V(x) = \begin{cases} 0, & |x| \le a \quad \text{(Region II: Well interior)} \\ V_0, & |x| > a \quad \text{(Regions I and III: Barriers)} \end{cases}
\end{equation}
For bound states ($0 < E < V_0$), let $k = \frac{\sqrt{2mE}}{\hbar}$ and $\alpha = \frac{\sqrt{2m(V_0 - E)}}{\hbar}$.

\textbf{2. Wavefunctions by Symmetry}:
\begin{itemize}
\item \textbf{Symmetric States (Even Parity):}
\begin{equation}
\psi_{II}(x) = B \cos(kx),\quad \psi_{III}(x) = C e^{-\alpha x}\; (x > a)
\end{equation}
Matching $\psi$ and $\frac{d\psi}{dx}$ at $x = a$:
\begin{equation}
B\cos(ka) = C e^{-\alpha a},\quad -kB\sin(ka) = -\alpha C e^{-\alpha a} \implies \mathbf{k \tan(ka) = \alpha}
\end{equation}
\item \textbf{Antisymmetric States (Odd Parity):}
\begin{equation}
\psi_{II}(x) = A \sin(kx) \implies \mathbf{-k \cot(ka) = \alpha}
\end{equation}
\end{itemize}

\textbf{3. Transcendental Graphical Solution}:
Let $\xi = ka$ and $\eta = \alpha a$.
\begin{equation}
\xi^2 + \eta^2 = \frac{2m V_0 a^2}{\hbar^2} \equiv R^2
\end{equation}
The allowed energy levels correspond to intersections of the circle of radius $R = \frac{\sqrt{2mV_0}a}{\hbar}$ with curves $\eta = \xi \tan\xi$ (even) and $\eta = -\xi \cot\xi$ (odd).
\begin{itemize}
\item At least one even bound state always exists no matter how shallow or narrow the well is.
\item Total number of bound states is finite: $N = 1 + \left\lfloor \frac{R}{\pi/2} \right\rfloor$.
\end{itemize}

\end{theorysol}
\vspace{4mm}

\begin{theoryq}{15}
Derive the reflection coefficient $R$ and transmission coefficient $T$ for a particle of energy $E$ incident on a potential step of height $V_0$ for both cases $E > V_0$ and $E < V_0$. Verify that $R + T = 1$.
\end{theoryq}
\begin{theorysol}
\textbf{1. Case I: Energy $E > V_0$}:
Let $k_1 = \frac{\sqrt{2mE}}{\hbar}$ for $x < 0$ and $k_2 = \frac{\sqrt{2m(E - V_0)}}{\hbar}$ for $x > 0$.
\begin{equation}
\psi_1(x) = A e^{j k_1 x} + B e^{-j k_1 x},\quad \psi_2(x) = C e^{j k_2 x}
\end{equation}
Applying boundary conditions at $x=0$ ($\psi_1(0)=\psi_2(0)$ and $\psi_1'(0)=\psi_2'(0)$):
\begin{equation}
A + B = C,\quad j k_1(A - B) = j k_2 C \implies k_1(A - B) = k_2(A + B)
\end{equation}
Solving for reflection amplitude $B/A$ and transmission amplitude $C/A$:
\begin{equation}
\frac{B}{A} = \frac{k_1 - k_2}{k_1 + k_2},\quad \frac{C}{A} = \frac{2k_1}{k_1 + k_2}
\end{equation}
The reflection and transmission coefficients are:
\begin{equation}
R = \left|\frac{B}{A}\right|^2 = \left(\frac{k_1 - k_2}{k_1 + k_2}\right)^2,\quad T = \frac{k_2}{k_1}\left|\frac{C}{A}\right|^2 = \frac{k_2}{k_1}\left(\frac{2k_1}{k_1 + k_2}\right)^2 = \frac{4k_1 k_2}{(k_1 + k_2)^2}
\end{equation}
\textbf{Conservation}:
\begin{equation}
R + T = \frac{(k_1 - k_2)^2 + 4k_1 k_2}{(k_1 + k_2)^2} = \frac{(k_1 + k_2)^2}{(k_1 + k_2)^2} = \mathbf{1}
\end{equation}

\textbf{2. Case II: Energy $E < V_0$}:
Here $k_2 = j\alpha$ where $\alpha = \frac{\sqrt{2m(V_0 - E)}}{\hbar}$.
\begin{equation}
\frac{B}{A} = \frac{k_1 - j\alpha}{k_1 + j\alpha} \implies R = \left|\frac{k_1 - j\alpha}{k_1 + j\alpha}\right|^2 = \frac{k_1^2 + \alpha^2}{k_1^2 + \alpha^2} = \mathbf{1},\quad T = 0
\end{equation}
The particle is totally reflected, but exhibits an exponentially decaying evanescent wave inside the barrier ($x > 0$).
\end{theorysol}
\vspace{4mm}

\begin{theoryq}{16}
Derive the Quantum Mechanical Probability Current Density $\mathbf{J}$ and prove the Probability Conservation Law (Continuity Equation) from the Time-Dependent Schrödinger Equation.
\end{theoryq}
\begin{theorysol}
\textbf{1. Time-Dependent Schrödinger Equation (TDSE)}:
\begin{equation}
j\hbar \frac{\partial \psi}{\partial t} = -\frac{\hbar^2}{2m}\nabla^2\psi + V(\mathbf{r})\psi \implies \frac{\partial \psi}{\partial t} = \frac{j\hbar}{2m}\nabla^2\psi - \frac{j}{\hbar}V\psi
\end{equation}
Taking the complex conjugate (assuming real potential $V^* = V$):
\begin{equation}
\frac{\partial \psi^*}{\partial t} = -\frac{j\hbar}{2m}\nabla^2\psi^* + \frac{j}{\hbar}V\psi^*
\end{equation}

\textbf{2. Time Derivative of Probability Density $\rho = \psi^* \psi$}:
\begin{align*}
\frac{\partial \rho}{\partial t} &= \frac{\partial}{\partial t}(\psi^* \psi) = \psi^* \frac{\partial \psi}{\partial t} + \frac{\partial \psi^*}{\partial t}\psi \\
&= \psi^* \left(\frac{j\hbar}{2m}\nabla^2\psi - \frac{j}{\hbar}V\psi\right) + \left(-\frac{j\hbar}{2m}\nabla^2\psi^* + \frac{j}{\hbar}V\psi^*\right)\psi \\
&= \frac{j\hbar}{2m} \left[\psi^* \nabla^2\psi - (\nabla^2\psi^*)\psi\right] = \frac{j\hbar}{2m} \nabla \cdot \left[\psi^* \nabla\psi - (\nabla\psi^*)\psi\right]
\end{align*}

\textbf{3. Defining Probability Current Density $\mathbf{J}$}:
\begin{equation}
\mathbf{J} = \frac{\hbar}{2mj} \left(\psi^* \nabla\psi - \psi \nabla\psi^*\right) = \frac{\hbar}{m}\text{Im}(\psi^* \nabla\psi)
\end{equation}
Substituting back yields the fundamental continuity equation:
\begin{equation}
\mathbf{\frac{\partial \rho}{\partial t} + \nabla \cdot \mathbf{J} = 0}
\end{equation}
Integrating over all space reveals that total probability is strictly conserved in time:
\begin{equation}
\frac{d}{dt}\int \rho\, dV = -\oint_{S_\infty} \mathbf{J} \cdot d\mathbf{S} = 0 \implies \int_{-\infty}^\infty |\psi(\mathbf{r},t)|^2 dV = 1\quad \forall t
\end{equation}
\end{theorysol}
\vspace{4mm}

\begin{theoryq}{17}
Construct a 1D Gaussian Wave Packet. Show that it represents a minimum uncertainty state ($\Delta x \Delta p = \hbar/2$) and derive the rate of wave packet spreading in time.
\end{theoryq}
\begin{theorysol}
\textbf{1. Gaussian Wave Packet Formulation}:
At $t=0$, let the spatial wavefunction be:
\begin{equation}
\psi(x,0) = \left(\frac{1}{2\pi \sigma_0^2}\right)^{1/4} e^{-\frac{x^2}{4\sigma_0^2}} e^{j k_0 x}
\end{equation}
The momentum probability amplitude $\phi(p)$ is the Fourier transform:
\begin{equation}
\phi(p) = \frac{1}{\sqrt{2\pi\hbar}}\int_{-\infty}^\infty \psi(x,0) e^{-j px/\hbar} dx = \left(\frac{2\sigma_0^2}{\pi\hbar^2}\right)^{1/4} e^{-\frac{\sigma_0^2 (p - p_0)^2}{\hbar^2}}
\end{equation}

\textbf{2. Calculating Uncertainties $\Delta x$ and $\Delta p$}:
\begin{equation}
\langle x \rangle = 0,\quad \langle x^2 \rangle = \sigma_0^2 \implies \Delta x = \sigma_0
\end{equation}
\begin{equation}
\langle p \rangle = \hbar k_0 = p_0,\quad \langle p^2 \rangle = p_0^2 + \frac{\hbar^2}{4\sigma_0^2} \implies \Delta p = \frac{\hbar}{2\sigma_0}
\end{equation}
The uncertainty product at $t=0$ is:
\begin{equation}
\mathbf{\Delta x \Delta p = \sigma_0 \left(\frac{\hbar}{2\sigma_0}\right) = \frac{\hbar}{2}}
\end{equation}
Thus, a Gaussian wave packet achieves the theoretical minimum uncertainty limit permitted by quantum mechanics.

\textbf{3. Time Evolution and Wave Packet Spreading}:
Solving TDSE for a free particle ($V=0$):
\begin{equation}
\sigma(t) = \sigma_0 \sqrt{1 + \left(\frac{\hbar t}{2m \sigma_0^2}\right)^2}
\end{equation}
As time $t$ advances, phase velocities of component Fourier waves differ ($v_p = \hbar k/2m$), causing the wave packet width to broaden relentlessly in coordinate space.
\end{theorysol}
\vspace{4mm}

\begin{theoryq}{18}
Explain the operating principle of the Scanning Tunneling Microscope (STM). Derive the exponential sensitivity of tunneling current to tip-to-sample distance $d$ and discuss constant-current vs constant-height imaging modes.
\end{theoryq}
\begin{theorysol}
\textbf{1. Quantum Tunneling Basis of STM}:
In classical mechanics, an electron cannot cross an insulating vacuum gap between a sharp conducting tip and a sample surface if its kinetic energy is less than the work function $\Phi$.
In quantum mechanics, the electron wavefunction penetrates the vacuum barrier of width $d$ with decay constant:
\begin{equation}
\kappa = \frac{\sqrt{2m_e \Phi}}{\hbar} \approx 0.512\sqrt{\Phi\text{ [eV]}}\; \text{\AA}^{-1}
\end{equation}

\textbf{2. Tunneling Current Dependence on Distance $d$}:
For a small bias voltage $V_b \ll \Phi$, the tunneling current $I_T$ is proportional to the barrier transmission probability $T$:
\begin{equation}
\mathbf{I_T \propto V_b e^{-2\kappa d}}
\end{equation}
For typical work functions $\Phi \approx 4.0\text{ eV}$, $\kappa \approx 1.0\text{ \AA}^{-1} = 10\text{ nm}^{-1}$:
\begin{equation}
2\kappa \approx 2.0\text{ \AA}^{-1} = 20\text{ nm}^{-1}
\end{equation}
Changing the tip-sample distance $d$ by merely $1\text{ \AA} = 0.1\text{ nm}$ alters the tunneling current by a factor of $e^{2} \approx 7.4$ (nearly an order of magnitude). This extreme sensitivity enables sub-Angstrom ($0.01\text{ nm}$) vertical resolution.

\textbf{3. Imaging Modes}:
\begin{itemize}
\item \textbf{Constant-Current Mode:} A feedback loop adjusts piezoelectric $z$-position to maintain constant $I_T$; recorded $z(x,y)$ voltage generates real-space atomic topography.
\item \textbf{Constant-Height Mode:} The tip scans at fixed $z$; rapid variations in $I_T(x,y)$ are recorded, enabling high scan speeds on atomically flat surfaces.
\end{itemize}

\end{theorysol}
\vspace{4mm}

\begin{theoryq}{19}
Derive the Energy-Time Uncertainty Principle $\Delta E \Delta t \ge \hbar/2$. Explain its physical interpretation using the lifetime $\tau$ of excited atomic states and the natural spectral line broadening $\Delta \nu$.
\end{theoryq}
\begin{theorysol}
\textbf{1. Derivation via Generalized Uncertainty Relation}:
For any observable operator $\hat{A}$ that does not explicitly depend on time:
\begin{equation}
\frac{d\langle \hat{A}\rangle}{dt} = \frac{1}{j\hbar}\langle [\hat{A}, \hat{H}] \rangle
\end{equation}
Using the Robertson-Schrödinger uncertainty inequality $\Delta A \Delta H \ge \frac{1}{2}|\langle [\hat{A}, \hat{H}] \rangle|$:
\begin{equation}
\Delta A \Delta E \ge \frac{\hbar}{2}\left|\frac{d\langle \hat{A}\rangle}{dt}\right|
\end{equation}
Defining the characteristic time $\Delta t$ required for the expectation value of $\hat{A}$ to change by one standard deviation:
\begin{equation}
\Delta t \equiv \frac{\Delta A}{\left|\frac{d\langle \hat{A}\rangle}{dt}\right|}
\end{equation}
Substituting gives:
\begin{equation}
\mathbf{\Delta E \Delta t \ge \frac{\hbar}{2}}
\end{equation}

\textbf{2. Physical Significance in Atomic Transitions}:
If an atom is excited to a state with finite lifetime $\tau = \Delta t$, its energy cannot be strictly sharply defined. The energy uncertainty is:
\begin{equation}
\Delta E \approx \frac{\hbar}{\tau} = \frac{h}{2\pi \tau}
\end{equation}
Since $\Delta E = h \Delta\nu$, the natural spectral linewidth is:
\begin{equation}
\Delta\nu_{nat} = \frac{1}{2\pi \tau}
\end{equation}
For typical atomic optical lifetimes $\tau \sim 10\text{ ns} = 10^{-8}\text{ s}$, the unavoidable natural spectral broadening is $\Delta\nu_{nat} \approx \frac{1}{2\pi(10^{-8})} \approx 16\text{ MHz}$.
\end{theorysol}
\vspace{4mm}

\begin{theoryq}{20}
An electron is trapped in an infinite potential well of width $L = 1.0\text{ nm}$. (a) Calculate the first three energy levels in eV. (b) Find the wavelength of the emitted photon when the electron falls from $n=4$ to $n=2$. (c) Calculate the expectation value $\langle x \rangle$ and $\langle x^2 \rangle$ for the ground state, and determine the spatial uncertainty $\Delta x$.
\end{theoryq}
\begin{theorysol}
\textbf{1. Energy Levels $E_1, E_2, E_3$}:
\begin{align*}
E_1 &= \frac{h^2}{8 m_e L^2} = \frac{(6.626 \times 10^{-34})^2}{8(9.109 \times 10^{-31})(1.0 \times 10^{-9})^2} = \frac{4.3904 \times 10^{-67}}{7.2872 \times 10^{-48}} \\
&= 6.0248 \times 10^{-20}\text{ J} = \frac{6.0248 \times 10^{-20}}{1.602 \times 10^{-19}} \approx \mathbf{0.376\text{ eV}} \\
E_2 &= 4 E_1 = 4(0.376\text{ eV}) = \mathbf{1.504\text{ eV}} \\
E_3 &= 9 E_1 = 9(0.376\text{ eV}) = \mathbf{3.384\text{ eV}} \\
E_4 &= 16 E_1 = 16(0.376\text{ eV}) = \mathbf{6.016\text{ eV}}
\end{align*}

\textbf{2. Transition Wavelength ($n=4 \to n=2$)}:
\begin{align*}
\Delta E &= E_4 - E_2 = 6.016 - 1.504 = 4.512\text{ eV} = 7.228 \times 10^{-19}\text{ J} \\
\lambda &= \frac{hc}{\Delta E} = \frac{(6.626 \times 10^{-34})(3.0 \times 10^8)}{7.228 \times 10^{-19}} = \frac{1.9878 \times 10^{-25}}{7.228 \times 10^{-19}} = 2.75 \times 10^{-7}\text{ m} = \mathbf{275\text{ nm}}\quad (\text{UV})
\end{align*}

\textbf{3. Expectation Values for Ground State ($n=1$)}:
Wavefunction $\psi_1(x) = \sqrt{\frac{2}{L}}\sin\left(\frac{\pi x}{L}\right)$.
\begin{equation}
\mathbf{\langle x \rangle} = \int_0^L x |\psi_1(x)|^2 dx = \frac{2}{L}\int_0^L x \sin^2\left(\frac{\pi x}{L}\right)dx = \mathbf{\frac{L}{2} = 0.50\text{ nm}}
\end{equation}
\begin{equation}
\langle x^2 \rangle = \frac{2}{L}\int_0^L x^2 \sin^2\left(\frac{\pi x}{L}\right)dx = L^2\left(\frac{1}{3} - \frac{1}{2\pi^2}\right) = L^2(0.3333 - 0.05066) = \mathbf{0.2827 L^2}
\end{equation}
The spatial dispersion is:
\begin{equation}
\mathbf{\Delta x} = \sqrt{\langle x^2 \rangle - \langle x \rangle^2} = \sqrt{0.2827 L^2 - 0.2500 L^2} = \sqrt{0.0327 L^2} \approx 0.1807 L = \mathbf{0.181\text{ nm}}
\end{equation}
\end{theorysol}
\vspace{4mm}
